\documentclass[11pt,a4paper]{article}

% ── Packages ──────────────────────────────────────────────────────────────────
\usepackage[utf8]{inputenc}
\usepackage[T1]{fontenc}
\usepackage{lmodern}
\usepackage{amsmath,amssymb,amsthm}
\usepackage{mathtools}
\usepackage{algorithm}
\usepackage[noend]{algpseudocode}
\usepackage{hyperref}
\usepackage[margin=2.5cm]{geometry}
\usepackage{enumitem}
\usepackage{booktabs}
\usepackage{xcolor}
\usepackage{tikz}
\usetikzlibrary{arrows.meta,positioning,calc}

% ── Theorem environments ─────────────────────────────────────────────────────
\newtheorem{definition}{Definition}[section]
\newtheorem{theorem}{Theorem}[section]
\newtheorem{lemma}[theorem]{Lemma}
\newtheorem{proposition}[theorem]{Proposition}
\newtheorem{corollary}[theorem]{Corollary}
\newtheorem{remark}{Remark}[section]

% ── Custom commands ───────────────────────────────────────────────────────────
\newcommand{\F}{\mathbb{F}}
\newcommand{\Z}{\mathbb{Z}}
\newcommand{\G}{\mathbb{G}}
\newcommand{\Gbj}{\G_{\mathrm{BJJ}}}
\newcommand{\Gbn}{\G_{\mathrm{BN}}}
\newcommand{\Fr}{\F_r}
\newcommand{\Fl}{\Z_\ell}
\newcommand{\Fp}{\F_p}
\newcommand{\order}{\mathrm{ord}}
\newcommand{\cofactor}{h}
\newcommand{\gen}{G}
\newcommand{\genbjj}{G_{\mathrm{BJJ}}}
\newcommand{\genbn}{G_{\mathrm{BN}}}
\newcommand{\sk}{\mathsf{sk}}
\newcommand{\pk}{\mathsf{pk}}
\newcommand{\PK}{\mathsf{PK}}
\newcommand{\SK}{\mathsf{SK}}
\newcommand{\id}{\mathsf{id}}
\newcommand{\sid}{\mathsf{sid}}
\newcommand{\pad}{\mathsf{pad}}
\newcommand{\hash}{\mathcal{H}}
\newcommand{\hashcurve}{\mathsf{H}}
\newcommand{\poly}{f}
\newcommand{\Poly}{\mathsf{Poly}}
\newcommand{\VSS}{\mathsf{VSS}}
\newcommand{\Commit}{\mathsf{Commit}}
\newcommand{\Verify}{\mathsf{Verify}}
\newcommand{\Prove}{\mathsf{Prove}}
\newcommand{\Extract}{\mathsf{Extract}_X}
\newcommand{\DH}{\mathsf{DH}}
\newcommand{\LHL}{\mathsf{LHL}}
\newcommand{\dhx}{\chi}
\newcommand{\eVRF}{\mathsf{eVRF}}
\newcommand{\PLONK}{\mathsf{PLONK}}
\newcommand{\FROST}{\mathsf{FROST}}
\newcommand{\bigO}{\mathcal{O}}
\newcommand{\negl}{\mathsf{negl}}
\newcommand{\getsr}{\xleftarrow{\$}}
\DeclareMathOperator{\ScalarMul}{ScalarMul}

% ── Title ─────────────────────────────────────────────────────────────────────
\title{%
  \textbf{GOLDEN Non-Interactive DKG}\\[0.3em]
  \large Adaptation to BN254/Baby~Jubjub with\\
  gnark PLONK eVRF Proofs and FROST Integration
}
\author{Fy Cryptographic Library --- Technical Specification}
\date{February 2026}

\begin{document}
\maketitle

\begin{abstract}
We present the design and implementation of a non-interactive distributed key
generation (DKG) protocol based on the GOLDEN construction (IACR~2025/1924),
adapted from BLS12-381 to the BN254/Baby~Jubjub curve pairing. The protocol
operates across two elliptic curve groups simultaneously: Baby~Jubjub
(BJJ) for Diffie--Hellman key agreement and identity proofs, and BN254~G1 for
polynomial secret sharing, Feldman VSS commitments, and threshold key
generation. An evaluatable Verifiable Random Function (eVRF), instantiated
via a gnark PLONK zero-knowledge proof system, enables non-interactive
encrypted share distribution without secure channels. The resulting key shares
are directly compatible with the FROST threshold Schnorr signature scheme
(RFC~9591~\cite{rfc9591}). This document provides a complete mathematical specification of
every component: the algebraic setting, cross-field conversions, hash-to-curve
construction, Shamir secret sharing, Feldman VSS, Schnorr identity proofs, the
eVRF pad derivation with Leftover Hash Lemma combination, the PLONK arithmetic
circuit, the full DKG protocol, and the FROST signing integration.
\end{abstract}

\tableofcontents
\newpage

% ══════════════════════════════════════════════════════════════════════════════
\section{Introduction}
\label{sec:intro}

Distributed key generation (DKG) is a foundational primitive for threshold
cryptography: $n$~participants jointly generate a shared secret key such that
any subset of $t$~or more can reconstruct it (or use it for signing), while
any~$t-1$ or fewer learn nothing about the secret. Traditional DKG protocols
require multiple interactive rounds and secure channels between participants.

The GOLDEN protocol~\cite{golden2025} achieves \emph{non-interactive} DKG by
replacing secure channels with an eVRF-based encryption scheme. Each dealer
encrypts shares for recipients using pads derived from Diffie--Hellman shared
secrets, proved correct via zero-knowledge proofs. The entire dealing is
broadcast in a single message.

This document describes our adaptation of GOLDEN from BLS12-381 to the
BN254/Baby~Jubjub curve pairing. The adaptation exploits a crucial algebraic
property: the base field of Baby~Jubjub \emph{coincides} with the scalar field
of BN254, enabling zero-cost x-coordinate extraction for the eVRF
construction. The output key shares integrate directly with the FROST threshold
Schnorr signature scheme.

\subsection{Notation}

Throughout this document, we use the following conventions:

\begin{center}
\renewcommand{\arraystretch}{1.2}
\begin{tabular}{@{}ll@{}}
\toprule
\textbf{Symbol} & \textbf{Meaning} \\
\midrule
$\getsr$ & Sampled uniformly at random \\
$[s]P$ & Scalar multiplication of point $P$ by scalar $s$ \\
$\mathcal{O}$ & The identity point (point at infinity) \\
$n$ & Total number of participants \\
$t$ & Signing threshold ($t \ge 2$) \\
$\sid$ & Session identifier ($\{0,1\}^{256}$) \\
$i \in [1, n]$ & Participant index (node ID) \\
$\|$ & Concatenation \\
\bottomrule
\end{tabular}
\end{center}

\subsection{Security Model}
\label{sec:secmodel-intro}

The GOLDEN DKG is analyzed in the following security model:

\begin{itemize}[nosep]
  \item \textbf{Static adversary.} A computationally bounded adversary corrupts
    up to $t - 1$ participants \emph{before} the protocol begins. Corrupted
    parties may deviate arbitrarily from the protocol (Byzantine faults). The
    adversary is \emph{rushing}: it may wait to see all honest messages before
    submitting its own.

  \item \textbf{No adaptive corruption.} The adversary's corruption set is fixed
    before Round~0. Adaptive corruption (choosing whom to corrupt based on
    protocol messages) is \emph{not} considered.

  \item \textbf{Authenticated broadcast with eventual delivery.} All Round0Msg
    messages are delivered authentically and eventually to all participants. The
    Schnorr identity proof (Section~\ref{sec:schnorr}) binds each dealing to
    the dealer's BJJ public key; transport-layer integrity is assumed.

  \item \textbf{Computational assumptions:}
    \begin{itemize}[nosep]
      \item Decisional Diffie--Hellman (DDH) on $\Gbj$ (Baby Jubjub).
      \item Discrete Logarithm (DL) on $\Gbn$ (BN254 G1).
      \item PLONK knowledge soundness (under the algebraic group model and
        random oracle model, as per~\cite{gabizon2019}).
    \end{itemize}
\end{itemize}

% ══════════════════════════════════════════════════════════════════════════════
\section{Algebraic Setting}
\label{sec:algebra}

The protocol operates over two distinct elliptic curve groups simultaneously.
The interplay between their field structures is the key enabler for the eVRF
construction.

% ── 2.1 ──────────────────────────────────────────────────────────────────────
\subsection{Baby Jubjub ($\Gbj$)}
\label{sec:bjj}

Baby Jubjub (BJJ) is a twisted Edwards curve defined over $\Fr$, the BN254
scalar field (i.e., the BJJ base field \emph{is} $\Fr$; see
Section~\ref{sec:fieldrel}). The curve equation is:
\begin{equation}
  \label{eq:bjj}
  a x^2 + y^2 = 1 + d\, x^2 y^2, \qquad a = -1,
\end{equation}
where $d = 12181644023421730124874158521699555681764249180949974110617291017600649128846$ is the curve parameter from gnark-crypto (see~\cite{lefkovits2017}).

\begin{remark}
  Setting $a = -1$ and rearranging~\eqref{eq:bjj} yields
  $-x^2 + y^2 = 1 + d\,x^2 y^2$, hence
  $y^2(1 - d\,x^2) = 1 + x^2$, which gives the $y^2$ formula used
  in hash-to-curve (equation~\eqref{eq:ysquared}).
\end{remark}

The full curve group has
order $8\ell$, where:
\begin{align}
  \ell &= 2736030358979909402780800718157159386076813972158567259200215660948447373041
  \label{eq:ell}
\end{align}
is the prime subgroup order ($\approx 2^{251}$ bits;
$\ell = \texttt{0x060C89CE5C263405370A08B6D0302B0BAB3EEDB83920EE0A677297DC392126F1}$),
and the cofactor is $\cofactor = 8$. We work exclusively in the prime-order
subgroup of order $\ell$.

\begin{itemize}[nosep]
  \item \textbf{Scalar field:} $\Fl$ (integers modulo $\ell$, $\approx 2^{251}$ bits).
  \item \textbf{Base field:} $\Fr$ (the BN254 scalar field, integers modulo $r$, $\approx 2^{254}$ bits).
  \item \textbf{Generator:} $\genbjj$ is the standard BJJ base point from gnark-crypto.
\end{itemize}

\noindent
\textbf{Note on naming:} Some references denote the BJJ base field as ``$\Fp$''
because Baby Jubjub is defined over a prime field. In this document, we
consistently write $\Fr$ for this field because it is the BN254 scalar field;
see~\eqref{eq:fieldrel}. The macro $\Fp$ is retained only for the BN254 base
field $\Fp'$.

% ── 2.2 ──────────────────────────────────────────────────────────────────────
\subsection{BN254 G1 ($\Gbn$)}
\label{sec:bn254}

BN254 G1 (also known as alt\_bn128) is a prime-order elliptic curve group
defined over the BN254 base field $\Fp'$. It has prime order:
\begin{align}
  r &= 21888242871839275222246405745257275088548364400416034343698204186575808495617
  \label{eq:r}
\end{align}
($r = \texttt{0x30644E72E131A029B85045B68181585D2833E84879B9709143E1F593F0000001}$),
with cofactor $1$ (every point except the identity is a generator of the full
group).

\begin{itemize}[nosep]
  \item \textbf{Scalar field:} $\Fr$ (integers modulo $r$, $\approx 2^{254}$ bits).
  \item \textbf{Base field:} $\Fp'$ (the BN254 base field, $\approx 2^{254}$ bits, distinct from $\Fr$).
  \item \textbf{Generator:} $\genbn$ is the standard BN254 G1 generator.
\end{itemize}

\begin{remark}[Concrete Security Level]
  \label{rem:bn254-security}
  BN254 provides approximately $100$--$110$ bits of security against the best
  known discrete-logarithm attacks (tower NFS~\cite{kb2016}), \emph{not} the
  128-bit level commonly assumed for generic prime-order groups. This is
  sufficient for many deployed systems (Ethereum, Zcash Sapling), but
  applications requiring $\ge 128$-bit security should migrate to BLS12-381 or
  BN382. The GOLDEN DKG construction is curve-agnostic via the
  \texttt{CurveSuite} abstraction and can be instantiated with higher-security
  pairings without protocol changes.
\end{remark}

% ── 2.3 ──────────────────────────────────────────────────────────────────────
\subsection{The Critical Field Relationship}
\label{sec:fieldrel}

The central algebraic property exploited by this construction is:

\begin{equation}
  \boxed{\text{BJJ base field} = \Fr = \text{BN254 scalar field}}
  \label{eq:fieldrel}
\end{equation}

This has three immediate consequences:

\begin{enumerate}[nosep]
  \item \textbf{Native x-coordinate extraction.} The coordinates $(x, y)$ of
        any BJJ point are elements of $\Fr$. Extracting the x-coordinate of a
        BJJ point yields an $\Fr$ element with no field conversion or reduction
        required.

  \item \textbf{Injection $\Fl \hookrightarrow \Fr$.} Since $\ell < r$, every
        element of $\Fl$ is a valid element of $\Fr$. The map
        $\iota: \Fl \to \Fr$ defined by $\iota(s) = s$ (interpreting the
        integer representative in $[0, \ell)$ as an element of $\Fr$) is an
        injection.

  \item \textbf{Reduction $\Fr \twoheadrightarrow \Fl$.} The map
        $\pi: \Fr \to \Fl$ defined by $\pi(s) = s \bmod \ell$ is a surjection.
        For any point $H$ of order $\ell$ and any $s \in \Fr$:
        \begin{equation}
          [s]H = [s \bmod \ell]H
          \label{eq:reduction}
        \end{equation}
        This is because $[\ell]H = \mathcal{O}$, so the scalar multiplication
        depends only on $s \bmod \ell$.
\end{enumerate}

\begin{remark}
  The field identity~\eqref{eq:fieldrel} is the reason BJJ was chosen as the
  ``inner'' curve for the eVRF: it enables the PLONK circuit
  (Section~\ref{sec:circuit}) to perform BJJ operations natively (without
  emulated arithmetic), dramatically reducing circuit size.
\end{remark}

% ── 2.4 ──────────────────────────────────────────────────────────────────────
\subsection{Field Conversion Functions}
\label{sec:fieldconv}

We define three conversion functions used throughout the protocol:

\begin{definition}[X-Coordinate Extraction]
  \label{def:extractx}
  $\Extract: \Gbj \to \Fr$ maps a BJJ point $P = (x, y)$ to its x-coordinate
  $x \in \Fr$. Since BJJ coordinates are natively in $\Fr$
  (Section~\ref{sec:fieldrel}), this is an exact injection---no modular
  reduction is performed.
\end{definition}

\begin{definition}[BJJ-to-Fr Injection]
  \label{def:bjjtofr}
  $\iota: \Fl \to \Fr$ interprets a BJJ scalar as a BN254 scalar. Since
  $\ell < r$, this is a lossless injection: every $\Fl$ element has a unique
  representative in $[0, \ell) \subset [0, r)$.
\end{definition}

\begin{definition}[Fr-to-BJJ Reduction]
  \label{def:frtobjj}
  $\pi: \Fr \to \Fl$ reduces a BN254 scalar modulo $\ell$:
  $\pi(s) = s \bmod \ell$. This is lossy for values $s \ge \ell$, but
  algebraically equivalent for subgroup operations by~\eqref{eq:reduction}.
\end{definition}

% ══════════════════════════════════════════════════════════════════════════════
\section{Cryptographic Primitives}
\label{sec:primitives}

% ── 3.1 ──────────────────────────────────────────────────────────────────────
\subsection{Hash Functions}
\label{sec:hash}

\subsubsection{Hash-to-Scalar}

We define two hash-to-scalar functions, one for each scalar field:

\begin{definition}[Hash-to-Scalar]
  \label{def:hashtoscalar}
  For a field $\F_q$ (either $\Fr$ or $\Fl$), the function
  $\hash_q: \{0,1\}^* \to \F_q$
  is constructed via counter-based SHA-256 expansion.

  \medskip\noindent
  \textbf{Input encoding.} Each Group implementation prepends a
  \emph{group-level domain prefix} $\mathsf{gdom}$ as raw bytes
  (e.g., \texttt{"BN254G1"} for $\Fr$, \texttt{"CircomBJJ"} for $\Fl$).
  Each data element $d_i$ (including any protocol-level domain separator
  passed as $d_0$) is a byte string prefixed with its length as a 4-byte
  big-endian unsigned integer.
  Point arguments are encoded as their canonical compressed byte representation
  before length-prefixing.

  \medskip\noindent
  \textbf{Expansion.}
  \begin{align}
    h_0 &= \mathrm{SHA\text{-}256}(\mathsf{gdom} \,\|\, \mathrm{len}(d_0) \,\|\, d_0 \,\|\, \cdots \,\|\, \texttt{0x00}) \\
    h_1 &= \mathrm{SHA\text{-}256}(\mathsf{gdom} \,\|\, \mathrm{len}(d_0) \,\|\, d_0 \,\|\, \cdots \,\|\, \texttt{0x01}) \\
    s &= (h_0 \,\|\, h_1) \bmod q
  \end{align}
  The 64-byte expanded value is interpreted as a big-endian integer and reduced
  modulo $q$. The statistical bias is at most $2^{-128}$ for any $q < 2^{256}$.
\end{definition}

\subsubsection{Hash-to-Curve (Try-and-Increment)}
\label{sec:hashtocurve}

Two independent hash-to-curve functions $\hashcurve_1, \hashcurve_2$ map
session data to BJJ subgroup points.

\begin{definition}[Hash-to-Curve]
  \label{def:hashtocurve}
  For domain $\mathsf{dom} \in \{\texttt{"golden-evrf-h1"}, \texttt{"golden-evrf-h2"}\}$
  and input data $(d_0, d_1, \ldots)$:
  \begin{enumerate}[nosep]
    \item For $\mathsf{ctr} = 0, 1, \ldots, 255$:
    \begin{enumerate}[nosep]
      \item Compute
        $\mathsf{digest} = \mathrm{SHA\text{-}256}\bigl(\mathsf{dom} \,\|\, \mathrm{len}(d_0) \,\|\, d_0 \,\|\, \cdots \,\|\, \mathsf{ctr}\bigr)$.
      \item Set $x = \mathsf{digest} \bmod r$ (interpret as an element of $\Fr$,
            which is also the BJJ base field per~\eqref{eq:fieldrel}).
      \item Solve the twisted Edwards equation~\eqref{eq:bjj} for $y$:
        \begin{equation}
          y^2 = \frac{1 + x^2}{1 - d \cdot x^2}
          \label{eq:ysquared}
        \end{equation}
      \item If the denominator $1 - d \cdot x^2 = 0$, continue.
      \item If $y^2$ is not a quadratic residue in $\Fr$ (Legendre symbol $\ne 1$),
            continue.
      \item Compute $y = \sqrt{y^2}$; select the \emph{lexicographically smallest}
            root by comparing the big-endian byte representations of $y$ and $-y$.
      \item Verify $(x, y)$ is on the curve.
      \item \textbf{Cofactor clearing:} $P_8 = [8](x, y)$.
      \item If $P_8 = \mathcal{O}$, continue.
      \item Verify subgroup membership: $[\ell]P_8 = \mathcal{O}$.
      \item Return $P_8$.
    \end{enumerate}
    \item If all 256 attempts fail, abort with error.
  \end{enumerate}
\end{definition}

\begin{remark}[Timing Side-Channel]
  Try-and-increment is \emph{not} constant-time: the number of iterations leaks
  through timing. This is acceptable because the inputs
  $\mathsf{sessData} = (\sid, \mathsf{randomMsg})$ become public once the
  dealing is broadcast (the \texttt{randomMsg} is included in the
  \texttt{Round0Msg}). During dealing \emph{construction}, the
  \texttt{randomMsg} is known only to the dealer, so a local
  timing side-channel is harmless (the dealer already knows its own secret).
  Verifiers call $\hashcurve$ only after receiving the broadcast, at which
  point all inputs are public. Secret-dependent inputs must never be passed
  to $\hashcurve$.
\end{remark}

\begin{remark}
  The domain separation $\texttt{"golden-evrf-h1"} \ne \texttt{"golden-evrf-h2"}$
  ensures $\hashcurve_1$ and $\hashcurve_2$ produce independent outputs for
  identical inputs, which is required for the Leftover Hash Lemma application
  (Section~\ref{sec:lhl}).
\end{remark}

% ── 3.2 ──────────────────────────────────────────────────────────────────────
\subsection{Shamir Secret Sharing}
\label{sec:shamir}

\begin{definition}[Random Polynomial]
  \label{def:polynomial}
  Given a group $\G$ with scalar field $\F_q$, a secret
  $\omega \in \F_q$, and degree $d$, define the random polynomial:
  \begin{equation}
    f(x) = \omega + a_1 x + a_2 x^2 + \cdots + a_d x^d, \qquad a_i \getsr \F_q
  \end{equation}
  The constant term $f(0) = \omega$ is the shared secret.
\end{definition}

\begin{definition}[Share Generation]
  \label{def:shares}
  For $n$ participants with IDs $\{1, 2, \ldots, n\}$, the share of participant
  $i$ is:
  \begin{equation}
    s_i = f(i) \in \F_q
  \end{equation}
  Evaluation uses Horner's method: starting from $\mathsf{result} = a_d$, iterate
  $\mathsf{result} \leftarrow \mathsf{result} \cdot x + a_{i}$ for
  $i = d-1, \ldots, 0$.
\end{definition}

\begin{theorem}[Lagrange Reconstruction]
  \label{thm:lagrange}
  Given any $t = d + 1$ shares $\{(i_j, s_{i_j})\}_{j=1}^{t}$, the secret can
  be reconstructed as:
  \begin{equation}
    \omega = f(0) = \sum_{j=1}^{t} s_{i_j} \cdot \lambda_{i_j}
    \label{eq:lagrange}
  \end{equation}
  where the Lagrange coefficients are:
  \begin{equation}
    \lambda_{i_j} = \prod_{\substack{k=1 \\ k \ne j}}^{t}
      \frac{i_k}{i_k - i_j}
    \label{eq:lagrange-coeff}
  \end{equation}
  Any $t-1$ or fewer shares reveal no information about $\omega$
  (information-theoretic security).
\end{theorem}

\begin{remark}
  The equivalent formulation $\lambda_{i_j} = \prod_{k \ne j} \frac{i_k}{i_k - i_j}$
  is often used in implementations; the two forms are algebraically equal.
\end{remark}

% ── 3.3 ──────────────────────────────────────────────────────────────────────
\subsection{Feldman Verifiable Secret Sharing}
\label{sec:vss}

Feldman VSS~\cite{feldman1987} extends Shamir sharing with public
commitments that allow verification of share correctness without
revealing the shares.

\begin{definition}[VSS Commitment]
  \label{def:vsscommit}
  Given a polynomial $f(x) = \sum_{k=0}^{t-1} a_k x^k$ over $\Fr$ and a
  generator $\genbn \in \Gbn$, the VSS commitment vector is:
  \begin{equation}
    \mathbf{A} = (A_0, A_1, \ldots, A_{t-1}), \qquad A_k = [a_k]\genbn
  \end{equation}
  In particular, $A_0 = [\omega]\genbn$ is the commitment to the secret.
\end{definition}

\begin{definition}[Expected Share Commitment]
  \label{def:expectedshare}
  The expected commitment to the share $f(i)$ is:
  \begin{equation}
    \VSS(i) = \sum_{k=0}^{t-1} [i^k] A_k = [f(i)] \genbn
    \label{eq:vss-verify}
  \end{equation}
  This follows from the homomorphism of scalar multiplication:
  $\sum [i^k] [a_k] \genbn = [\sum a_k \cdot i^k] \genbn = [f(i)] \genbn$.
\end{definition}

\begin{remark}
  In this implementation, the polynomial $f$ is defined over $\Fr$ (the BN254
  scalar field) and commitments are computed in $\Gbn$. This ensures
  consistency: the Shamir shares, ciphertext arithmetic, and VSS verification
  all operate in the same field $\Fr$.
\end{remark}

% ── 3.4 ──────────────────────────────────────────────────────────────────────
\subsection{Schnorr Proof of Knowledge (Identity Proof)}
\label{sec:schnorr}

Each dealer proves knowledge of their BJJ secret key, bound to the session ID,
to prevent identity impersonation and cross-session replay.

\begin{definition}[Session-Bound Schnorr PoK]
  \label{def:schnorr}
  For a prover with secret key $\sk \in \Fl$ and public key
  $\pk = [\sk]\genbjj$, bound to session $\sid$:

  \medskip
  \textbf{Prove}$(\sk, \pk, \sid)$:
  \begin{enumerate}[nosep]
    \item Sample $\nu \getsr \Fl$ \hfill (nonce)
    \item $R \leftarrow [\nu]\genbjj$ \hfill (commitment)
    \item $c \leftarrow \hash_\ell\bigl(\texttt{"golden-pki-pok"} \,\|\, \sid \,\|\, \pk \,\|\, R\bigr)$
          \hfill (challenge)
    \item $s \leftarrow \nu - c \cdot \sk \pmod{\ell}$ \hfill (response)
    \item Securely erase $\nu$
    \item Return $\pi = (R, c, s)$
  \end{enumerate}

  \medskip
  \textbf{Verify}$(\pk, \sid, \pi = (R, c, s))$:
  \begin{enumerate}[nosep]
    \item $R' \leftarrow [s]\genbjj + [c]\pk$ \hfill (recompute commitment)
    \item $c' \leftarrow \hash_\ell\bigl(\texttt{"golden-pki-pok"} \,\|\, \sid \,\|\, \pk \,\|\, R'\bigr)$
    \item Accept if and only if $c' = c$
  \end{enumerate}
\end{definition}

\begin{proposition}[Correctness]
  If $\pi$ was honestly generated, then
  $R' = [s]\genbjj + [c]\pk = [\nu - c\sk]\genbjj + [c \cdot \sk]\genbjj = [\nu]\genbjj = R$,
  so $c' = c$.
\end{proposition}

\begin{remark}
  This uses the \emph{subtraction} convention ($s = \nu - c \cdot \sk$), with
  verification $R' = [s]G + [c]\pk$. The addition convention
  ($s = \nu + c \cdot \sk$ with $R' = [s]G - [c]\pk$) is equivalent and
  equally secure; the choice is purely a matter of implementation convention.
\end{remark}

% ── 3.5 ──────────────────────────────────────────────────────────────────────
\subsection{Domain Separation}
\label{sec:domains}

All hash function invocations use distinct domain separators to prevent
cross-protocol attacks:

\begin{center}
\renewcommand{\arraystretch}{1.3}
\begin{tabular}{@{}lllp{4cm}@{}}
\toprule
\textbf{Domain String} & \textbf{Type} & \textbf{Function} & \textbf{Purpose} \\
\midrule
\texttt{"BN254G1"} & group-level & $\hash_r$ & Raw prefix for all BN254 hash-to-scalar \\
\texttt{"CircomBJJ"} & group-level & $\hash_\ell$ & Raw prefix for all BJJ hash-to-scalar \\
\texttt{"golden-evrf-h1"} & protocol & $\hashcurve_1$ & First hash-to-curve for eVRF (raw) \\
\texttt{"golden-evrf-h2"} & protocol & $\hashcurve_2$ & Second hash-to-curve for eVRF (raw) \\
\texttt{"golden-lhl-alpha"} & protocol & $\hash_r$ & LHL coefficient (length-prefixed $d_0$) \\
\texttt{"golden-pki-pok"} & protocol & $\hash_\ell$ & Schnorr challenge (length-prefixed $d_0$) \\
\bottomrule
\end{tabular}
\end{center}

Group-level domain prefixes are written as raw bytes before the length-prefixed
data elements (see Definition~\ref{def:hashtoscalar}). Hash-to-curve functions
use their protocol domain string as a raw prefix directly (no length-prefixing).
Protocol-level domain strings passed to $\hash_q$ are length-prefixed as $d_0$,
preventing concatenation ambiguity attacks.

% ══════════════════════════════════════════════════════════════════════════════
\section{The eVRF Construction}
\label{sec:evrf}

The eVRF (evaluatable Verifiable Random Function) is the core cryptographic
mechanism enabling non-interactive share encryption. It derives a
deterministic, symmetric encryption pad from a Diffie--Hellman shared secret
over BJJ, combined with hash-to-curve evaluations via the Leftover Hash Lemma.

% ── 4.1 ──────────────────────────────────────────────────────────────────────
\subsection{Diffie--Hellman on BJJ}
\label{sec:dh}

Given a dealer with secret key $\sk_d \in \Fl$ and a recipient with public key
$\pk_j = [\sk_j]\genbjj$, the DH shared secret is:
\begin{equation}
  S_{d,j} = [\sk_d]\pk_j = [\sk_d \cdot \sk_j]\genbjj
  \label{eq:dh}
\end{equation}
This is symmetric: $[\sk_d]\pk_j = [\sk_j]\pk_d$ (both equal
$[\sk_d \sk_j]\genbjj$).

\begin{remark}[DerivePad Scalar Type]
  The $\mathsf{DerivePad}$ function accepts the dealer's secret key as a
  $\Fl$ scalar (an inner-group scalar in $\Z_\ell$). This is important because
  DH is performed on BJJ, where the group order is $\ell$, not $r$.
\end{remark}

\begin{remark}[Degenerate DH Check]
  If $S_{d,j} = \mathcal{O}$ (the identity point), the pad would be $0$,
  leaking the share in cleartext. The protocol aborts with
  $\mathsf{ErrDegenerateDH}$ if this occurs. Under honestly generated keys
  ($\sk \ne 0$), this probability is negligible.
\end{remark}

% ── 4.2 ──────────────────────────────────────────────────────────────────────
\subsection{X-Coordinate Extraction and Field Conversion}
\label{sec:xextract}

Write $S_{d,j} = (S_x, S_y)$ in BJJ affine coordinates. Since BJJ coordinates
are in $\Fr$ (Section~\ref{sec:fieldrel}):
\begin{equation}
  \dhx = \Extract(S_{d,j}) = S_x \in \Fr
  \label{eq:extract}
\end{equation}
To perform scalar multiplication on BJJ subgroup points, we need a scalar in
$\Fl$:
\begin{equation}
  \dhx_{\mathrm{BJJ}} = \pi(\dhx) = \dhx \bmod \ell \in \Fl
  \label{eq:sbjj}
\end{equation}
By~\eqref{eq:reduction}, for any point $H$ of order $\ell$:
$[\dhx]H = [\dhx_{\mathrm{BJJ}}]H$.

% ── 4.3 ──────────────────────────────────────────────────────────────────────
\subsection{VRF Evaluation}
\label{sec:vrfeval}

Let $\mathsf{sessData} = (\sid \,\|\, \mathsf{randomMsg}_d)$ be the session
data, where $\mathsf{randomMsg}_d \getsr \{0,1\}^{256}$ is a fresh nonce per
dealing (generated in Algorithm~\ref{alg:dealing}, line~7). Compute two
independent hash-to-curve evaluations:
\begin{align}
  H_1 &= \hashcurve_1(\mathsf{sessData}) \in \Gbj  \label{eq:h1} \\
  H_2 &= \hashcurve_2(\mathsf{sessData}) \in \Gbj  \label{eq:h2}
\end{align}
If $H_1 = \mathcal{O}$ or $H_2 = \mathcal{O}$, abort with
$\mathsf{ErrIdentityPoint}$. (This guards against degenerate hash-to-curve
outputs that would produce a zero pad regardless of the DH secret.)

\begin{remark}[Negligible Probability of $\dhx = 0$]
  \label{rem:chi-zero}
  The DH x-coordinate $\dhx = \Extract([\sk_d]\pk_j) \in \Fr$ equals zero with
  probability $1/r \approx 2^{-254}$, which is negligible. If $\dhx = 0$, then
  $P_1 = P_2 = \mathcal{O}$ and the pad is identically zero regardless of
  $\alpha$. The implementation does not explicitly check for this case; the
  $H_i = \mathcal{O}$ check above catches a different failure mode (degenerate
  hash-to-curve output). Since $\Pr[\dhx = 0] = \negl(\lambda)$, this is a
  non-issue in practice.
\end{remark}

Then evaluate the VRF at these points:
\begin{align}
  P_1 &= [\dhx_{\mathrm{BJJ}}] H_1 \in \Gbj  \label{eq:p1} \\
  P_2 &= [\dhx_{\mathrm{BJJ}}] H_2 \in \Gbj  \label{eq:p2}
\end{align}

% ── 4.4 ──────────────────────────────────────────────────────────────────────
\subsection{Randomness Extraction via the Computational Leftover Hash Lemma}
\label{sec:lhl}

The x-coordinate extraction function $\Extract$ maps BJJ points to $\Fr$ elements,
but it is not a random oracle---it is a deterministic algebraic 2-to-1 map
(on a twisted Edwards curve, $P$ and $-P$ share the same $x$-coordinate).
To extract computationally indistinguishable randomness from the DDH-hard
DH output, we use the \emph{computational} variant of the Leftover Hash
Lemma~\cite{dodis2008}.

\medskip\noindent
\textbf{Setup.} The classical (statistical) LHL~\cite{ill1989} requires the
source to have high \emph{min-entropy}: $H_\infty(X) \ge \log|\Fr| + 2\lambda$,
where $\lambda$ is the security parameter. Because $\Extract$ is 2-to-1,
each x-coordinate $\Extract(P_i)$ has at most $\log\ell - 1$ bits of
min-entropy (where $\ell$ is the BJJ subgroup order), which is
comfortably above $\log r$ since $r < \ell$. However, under DDH the source
is only \emph{computationally} hard, not information-theoretically uniform.
We therefore invoke the computational LHL of Dodis et al.~\cite{dodis2008}:

\begin{lemma}[Computational Leftover Hash Lemma {\cite[Lemma~2.4]{dodis2008}}]
  \label{lem:clhl}
  Let $\mathcal{H} = \{h_\alpha : \{0,1\}^n \to \{0,1\}^m\}_{\alpha \in \mathcal{K}}$
  be a family of universal hash functions, and let $X$ be a random variable
  over $\{0,1\}^n$ with \emph{computational} min-entropy
  $H_\infty^{\mathrm{HILL}}(X) \ge m + 2\lambda$. Then
  $(\alpha,\, h_\alpha(X))$ is computationally indistinguishable from
  $(\alpha,\, U_m)$ where $U_m$ is uniform over $\{0,1\}^m$ and $\alpha
  \getsr \mathcal{K}$.
\end{lemma}

\noindent\textbf{Application to GOLDEN.}
We identify $X = (\Extract(P_1),\, \Extract(P_2)) \in \Fr^2$ and
$h_\alpha(x_1, x_2) = x_1 + \alpha \cdot x_2 \in \Fr$ (an inner-product
hash, which is a universal hash family over $\Fr$). Under DDH on $\Gbj$,
the pair $(P_1, P_2) = ([\dhx]H_1,\, [\dhx]H_2)$ is computationally
indistinguishable from two independent random subgroup points. Since
$\Extract$ is 2-to-1, the joint source $X$ carries HILL min-entropy
at least $2(\log\ell - 1) \ge 2 \times 252 > 254 + 2\lambda$ bits
(for $\lambda = 100$, the BN254 security level; Remark~\ref{rem:bn254-security}).
By Lemma~\ref{lem:clhl}, $(\alpha,\, \pad)$ is computationally
indistinguishable from $(\alpha,\, U)$ where $U \getsr \Fr$.

\begin{definition}[LHL Pad Derivation]
  \label{def:pad}
  Given a public session-level coefficient
  \begin{equation}
    \alpha = \hash_r\bigl(\texttt{"golden-lhl-alpha"} \,\|\, \sid\bigr) \in \Fr
    \label{eq:alpha}
  \end{equation}
  the encryption pad is:
  \begin{equation}
    \pad = \Extract(P_1) + \alpha \cdot \Extract(P_2) \in \Fr
    \label{eq:pad}
  \end{equation}
  Expanding: $\pad = P_1.x + \alpha \cdot P_2.x$, where all arithmetic is in
  $\Fr$.
\end{definition}

\begin{remark}
  The coefficient $\alpha$ is derived from the session ID alone (not
  per-dealer or per-recipient). This is consistent with the GOLDEN protocol
  specification, where $\alpha$ is a session-level parameter. Since different
  dealer--recipient pairs have different DH secrets $S_{d,j}$, the pads are
  already unique per pair. By Lemma~\ref{lem:clhl}, the pad is computationally
  indistinguishable from a uniform $\Fr$ element given $\alpha$.
\end{remark}

% ── 4.5 ──────────────────────────────────────────────────────────────────────
\subsection{Pad Commitment}
\label{sec:padcommit}

The dealer commits to the pad by computing:
\begin{equation}
  R_{d,j} = [\pad]\genbn \in \Gbn
  \label{eq:rcommitment}
\end{equation}
This commitment serves two purposes:
\begin{enumerate}[nosep]
  \item It enables ciphertext consistency verification (Proposition~\ref{prop:ct-verify}).
  \item It is a public input to the eVRF PLONK proof (Section~\ref{sec:circuit}).
\end{enumerate}

% ── 4.6 ──────────────────────────────────────────────────────────────────────
\subsection{Share Encryption}
\label{sec:encrypt}

The dealer encrypts the share $s_j = f_d(j)$ for recipient $j$ as:
\begin{equation}
  z_{d,j} = \pad + s_j \in \Fr
  \label{eq:ciphertext}
\end{equation}
The pair $(R_{d,j}, z_{d,j})$ forms the ciphertext. Decryption by recipient
$j$ (who can re-derive $\pad$ via their own DH computation) recovers:
\begin{equation}
  s_j = z_{d,j} - \pad
  \label{eq:decrypt}
\end{equation}

% ── 4.7 ──────────────────────────────────────────────────────────────────────
\subsection{Symmetry Property}
\label{sec:symmetry}

The pad is symmetric between dealer $d$ and recipient $j$:
\begin{equation}
  \mathsf{DerivePad}(\sk_d, \pk_j, \mathsf{sessData}, \alpha)
  = \mathsf{DerivePad}(\sk_j, \pk_d, \mathsf{sessData}, \alpha)
  \label{eq:symmetry}
\end{equation}
This follows from the symmetry of Diffie--Hellman: $[\sk_d]\pk_j = [\sk_j]\pk_d$,
which produces the same shared secret $S$ and thus the same $\dhx, \dhx_{\mathrm{BJJ}}, P_1, P_2, \pad$.

% ══════════════════════════════════════════════════════════════════════════════
\section{eVRF Zero-Knowledge Proof}
\label{sec:zkproof}

The dealer must prove that the pad commitment $R_{d,j}$ was correctly derived
from the DH shared secret, without revealing their secret key $\sk_d$. This is
accomplished via a PLONK~\cite{gabizon2019} zero-knowledge proof over a
gnark~\cite{gnark} arithmetic circuit.

% ── 5.1 ──────────────────────────────────────────────────────────────────────
\subsection{Proof Statement}
\label{sec:proofstatement}

The eVRF proof establishes the following relation:

\begin{equation}
  \boxed{
    \exists\, \sk \in \Fl : \quad
    \begin{aligned}
      &\pk_d = [\sk]\genbjj \\
      &S = [\sk]\pk_j \\
      &\pad = \Extract([\dhx]H_1) + \alpha \cdot \Extract([\dhx]H_2) \\
      &R = [\pad]\genbn
    \end{aligned}
  }
  \quad\text{where } \dhx = \Extract(S)
  \label{eq:proof-relation}
\end{equation}

where $(\pk_d, \pk_j, H_1, H_2, \alpha, R)$ are public inputs and $\sk$ is
the private witness.

\begin{remark}[Witness Domain]
  The secret key $\sk$ lives in $\Fl$ (the BJJ scalar field), but the PLONK
  circuit is defined over $\Fr$. The circuit witness is $\iota(\sk) \in \Fr$
  (the injected value; see Definition~\ref{def:bjjtofr}). Since $\ell < r$,
  every $\Fl$ element has a unique representative in $\Fr$, and the BJJ
  scalar multiplication gadget in gnark correctly handles $\Fr$ witnesses
  operating on order-$\ell$ subgroup points (because $[\iota(\sk)]P = [\sk]P$
  for any $P$ of order $\ell$, by~\eqref{eq:reduction}).
\end{remark}

% ── 5.2 ──────────────────────────────────────────────────────────────────────
\subsection{Circuit Architecture}
\label{sec:circuit}

The PLONK arithmetic circuit is compiled over $\Fr$ (the BN254 scalar field),
which is the native field of the gnark constraint system for BN254-based
proofs. This choice has a crucial efficiency implication:

\begin{center}
\renewcommand{\arraystretch}{1.3}
\begin{tabular}{@{}lll@{}}
\toprule
\textbf{Operation} & \textbf{Curve} & \textbf{Cost} \\
\midrule
BJJ scalar mult. & $\Gbj$ & \emph{Native} (BJJ base field $=$ circuit field $\Fr$) \\
BJJ point add. & $\Gbj$ & \emph{Native} \\
BN254 G1 scalar mult. & $\Gbn$ & \emph{Emulated} (G1 base field $\Fp' \ne \Fr$) \\
\bottomrule
\end{tabular}
\end{center}

\subsubsection{Public Inputs}

The circuit has the following public input variables:
\begin{align}
  &\mathsf{DealerPK}_X, \mathsf{DealerPK}_Y \in \Fr
    &&\text{(BJJ point, native)} \label{eq:pub-dealerpk} \\
  &\mathsf{RecipientPK}_X, \mathsf{RecipientPK}_Y \in \Fr
    &&\text{(BJJ point, native)} \label{eq:pub-recipientpk} \\
  &H_{1,X}, H_{1,Y} \in \Fr
    &&\text{(BJJ point, native)} \label{eq:pub-h1} \\
  &H_{2,X}, H_{2,Y} \in \Fr
    &&\text{(BJJ point, native)} \label{eq:pub-h2} \\
  &\alpha \in \Fr
    &&\text{(scalar, native)} \label{eq:pub-alpha} \\
  &R_X, R_Y \in \Fp'
    &&\text{(BN254 G1 point, emulated)} \label{eq:pub-r}
\end{align}

\subsubsection{Private Witness}

\begin{equation}
  \mathsf{DealerSK} = \iota(\sk) \in \Fr \qquad \text{(the BJJ secret key, injected into $\Fr$)}
  \label{eq:priv-sk}
\end{equation}

\subsubsection{Constraint Steps}

\begin{enumerate}
  \item \textbf{Public Key Ownership} (BJJ, native):
  \begin{align}
    \mathsf{computedPK} &= [\mathsf{DealerSK}]\genbjj \\
    \mathsf{AssertOnCurve}(\mathsf{DealerPK}) & \\
    \mathsf{computedPK}.X &\stackrel{!}{=} \mathsf{DealerPK}_X \\
    \mathsf{computedPK}.Y &\stackrel{!}{=} \mathsf{DealerPK}_Y
  \end{align}

  \item \textbf{DH Shared Secret} (BJJ, native):
  \begin{align}
    \mathsf{AssertOnCurve}(\mathsf{RecipientPK}) & \\
    S &= [\mathsf{DealerSK}]\mathsf{RecipientPK}
  \end{align}

  \item \textbf{X-Coordinate Extraction} (native):
  \begin{equation}
    \dhx = S.X \in \Fr
  \end{equation}

  \begin{remark}[In-Circuit vs.\ Out-of-Circuit Equivalence]
    Out-of-circuit, the code reduces $\dhx \bmod \ell$ before BJJ scalar
    multiplication. In-circuit, $\dhx$ is used directly as a native $\Fr$
    variable. These are algebraically equivalent because $H_1, H_2$ are in the
    prime-order subgroup of order $\ell$, so $[\dhx]H = [\dhx \bmod \ell]H$.
    The gnark twisted~Edwards \texttt{ScalarMul} gadget correctly handles
    scalars in $\Fr$ operating on order-$\ell$ subgroup points.
  \end{remark}

  \item \textbf{Hash-to-Curve Multiplications} (BJJ, native):
  \begin{align}
    \mathsf{AssertOnCurve}(H_1),\; \mathsf{AssertOnCurve}(H_2) & \\
    P_1 &= [\dhx] H_1 \\
    P_2 &= [\dhx] H_2
  \end{align}

  \item \textbf{LHL Combination} (native $\Fr$ arithmetic):
  \begin{equation}
    \pad = P_1.X + \alpha \cdot P_2.X
  \end{equation}

  \item \textbf{BN254 G1 Commitment Verification} (emulated):
  \begin{align}
    \mathsf{padBits} &= \mathsf{ToBinary}(\pad, 254)
      &&\text{(decompose native $\Fr$ to 254 bits)} \\
    \mathsf{padScalar} &= \mathsf{FromBits}(\mathsf{padBits})
      &&\text{(reconstruct as emulated $\Fr$)} \\
    \mathsf{computedR} &= [\mathsf{padScalar}]\genbn
      &&\text{(emulated BN254 G1 scalar mult.)} \\
    \mathsf{computedR} &\stackrel{!}{=} R
  \end{align}
\end{enumerate}

\begin{remark}[Bit Decomposition Bridge]
  Step~6 is the most constraint-expensive part of the circuit. The pad $\pad$
  is a native $\Fr$ variable, but the emulated BN254 G1 arithmetic
  (\texttt{sw\_emulated}) expects an emulated $\Fr$ element with 4 limbs of 64
  bits. The bridge is a 254-bit decomposition: the native value is split into
  individual bits via range-check constraints, then reassembled as an emulated
  element. This costs approximately $254$ range-check constraints plus the
  emulated scalar multiplication ($\approx 10^4$ constraints).
\end{remark}

% ── 5.3 ──────────────────────────────────────────────────────────────────────
\subsection{PLONK Proof System}
\label{sec:plonk}

The circuit is compiled and proved using the PLONK proof system with KZG
polynomial commitments over BN254.

\begin{itemize}[nosep]
  \item \textbf{Circuit compilation:} The circuit is compiled once (lazily,
    via \texttt{sync.Once}) into a constraint system (SCS format).
  \item \textbf{Key generation:} Proving key $\mathsf{pk}$ and verifying key
    $\mathsf{vk}$ are generated from the SCS and a KZG structured reference
    string (SRS).
  \item \textbf{Proof generation:}
    $\mathsf{proof} \leftarrow \PLONK.\Prove(\mathsf{ccs}, \mathsf{pk}, \mathsf{witness})$.
  \item \textbf{Verification:}
    $\PLONK.\Verify(\mathsf{proof}, \mathsf{vk}, \mathsf{publicWitness})$.
\end{itemize}

\begin{remark}[SRS Security]
  \label{rem:srs}
  The production implementation uses the \textbf{Aztec Ignition ceremony SRS}~\cite{aztec-ignition},
  a BN254 powers-of-tau structured reference string produced by a multi-party
  computation with 176 participants. The ceremony transcript
  (\texttt{transcript00.dat}, ${\sim}307$~MB, containing ${\sim}5 \times 10^6$
  G1 points) is downloaded on first use and cached locally under
  \texttt{\$HOME/.cache/fy/srs/}. After parsing, a batch Schwartz--Zippel
  pairing verification~\cite{schwartz1980} checks consistency of all $\tau^i
  \cdot G_1$ powers against $\tau \cdot G_2$:
  \begin{equation}
    e\!\Bigl(\sum_i r^i \cdot [\tau^{i+1}]G_1,\; G_2\Bigr)
    = e\!\Bigl(\sum_i r^i \cdot [\tau^i]G_1,\; [\tau]G_2\Bigr)
    \label{eq:srs-verify}
  \end{equation}
  for a random challenge $r \getsr \Fr$. This catches download corruption,
  malformed transcripts, and inconsistent $G_1$/$G_2$ data in a single
  multi-pairing.

  For testing, a \texttt{testingSRSProvider} callback (set in
  \texttt{\_test.go} \texttt{init()}) overrides the ceremony SRS with
  \texttt{unsafekzg.NewSRS} (known toxic waste---\textbf{do not use in
  production}). A runtime guard panics if \texttt{testingSRSProvider} is set
  in a non-test binary.
\end{remark}

\begin{remark}[SRS Provenance vs.\ Integrity]
  \label{rem:srs-provenance}
  The pairing check~\eqref{eq:srs-verify} verifies \emph{internal consistency}
  (the $G_1$ powers form a geometric progression under a single $\tau$) but does
  \emph{not} verify \emph{provenance}: it cannot confirm that this particular
  $\tau$ was generated by the Ignition ceremony rather than by an adversary who
  knows $\tau$. Provenance assurance relies on (a)~downloading from the
  canonical Ignition verification repository, and (b)~optionally pinning the
  SRS to a known hash digest.  The current implementation verifies consistency
  but does not pin a hash; deployers requiring stronger provenance guarantees
  should add a SHA-256 digest check on the raw transcript before parsing.
\end{remark}

\begin{remark}[Lagrange SRS Cache Trust Model]
  \label{rem:lagrange-cache}
  After downloading and verifying the monomial SRS, the implementation precomputes
  a Lagrange-basis SRS (required by PLONK) and caches it to disk. This cached
  Lagrange SRS is loaded \emph{without re-verification} on subsequent runs, so
  the cache directory (\texttt{\$HOME/.cache/fy/srs/}) must be trusted.
  An adversary with write access to this directory could substitute a malicious
  Lagrange SRS, enabling forged proofs. Deployers should protect the cache with
  appropriate filesystem permissions, or delete the cache and re-derive on each
  startup if the deployment environment is untrusted.
\end{remark}

\subsubsection{Secret Key Zeroing}

After witness construction, the secret key material ($\sk$ as a
\texttt{big.Int} and its serialized \texttt{skBytes}) is securely erased by:
\begin{enumerate}[nosep]
  \item Obtaining the underlying word array via \texttt{big.Int.Bits()} (shared
    backing memory).
  \item Overwriting each word with zero.
  \item Calling \texttt{SetInt64(0)} to reset the logical value.
  \item Overwriting the serialized \texttt{skBytes} array element-by-element.
  \item Calling \texttt{runtime.KeepAlive} on both \texttt{skBigInt} and
    \texttt{skBytes} to prevent the compiler from optimizing away the zeroing
    by ensuring the variables are considered live until after the overwrite
    completes.
\end{enumerate}
Steps~1--3 are necessary because \texttt{SetInt64(0)} alone only sets the
length to zero without clearing the underlying memory. Step~5 is necessary
because the Go compiler may otherwise elide dead stores to variables that are
not read afterwards.

% ══════════════════════════════════════════════════════════════════════════════
\section{DKG Protocol}
\label{sec:dkg}

The GOLDEN DKG consists of three phases: dealing (broadcast), verification, and
completion (aggregation). It is non-interactive: each participant broadcasts a
single message containing encrypted shares for all peers.

% ── 6.1 ──────────────────────────────────────────────────────────────────────
\subsection{Parameters and Setup}
\label{sec:dkg-setup}

\begin{itemize}[nosep]
  \item $n$: total participants, $2 \le n \le 256$ (\texttt{MaxParticipants}).
    The upper bound is a DoS-prevention limit (the $\mathcal{O}(n^2)$ proof
    count becomes impractical for larger groups) rather than a security
    parameter.
  \item $t$: threshold, $2 \le t \le n$. The minimum threshold $t = 2$ is
    enforced for FROST compatibility: FROST's binding factor computation
    requires at least two signers. A threshold of $t = 1$ degenerates to
    single-party Schnorr signing.
  \item $\sid \in \{0,1\}^{256}$: unique session identifier.
  \item Each participant $i \in [1, n]$ holds $(\sk_i, \pk_i)$ where
    $\sk_i \getsr \Fl \setminus \{0\}$ and $\pk_i = [\sk_i]\genbjj$.
  \item All participants know all public keys $\{\pk_1, \ldots, \pk_n\}$.
  \item Node IDs must be in $[1, 2^{31} - 1]$ (\texttt{MaxNodeID}) to fit in
    \texttt{uint32} encoding for FROST Lagrange coefficient computation. The
    encoding places the ID as a big-endian \texttt{uint32} at bytes $[28{:}32]$
    of a 32-byte zero-padded buffer.
\end{itemize}

\begin{remark}[Subgroup Membership of Recipient Public Keys]
  \label{rem:subgroup-check}
  Baby Jubjub has cofactor $\cofactor = 8$. A malicious participant could
  register a public key $\pk^* \notin \Gbj$ (i.e., a point in the order-$8\ell$
  curve group but not in the prime-order subgroup of order~$\ell$). If
  \texttt{CreateDealing} uses such a $\pk^*$ in the DH computation, the
  resulting shared secret $[\sk_d]\pk^*$ may have low order, enabling a
  small-subgroup attack that leaks information about $\sk_d \bmod 8$.

  \medskip\noindent\textbf{Mitigation.}
  The implementation performs an identity check ($\pk_j \ne \mathcal{O}$)
  but does not currently enforce $[\ell]\pk_j = \mathcal{O}$. This is safe
  \emph{provided} that public keys are validated at registration time
  (e.g., by requiring a proof-of-possession or subgroup check). Deployments
  that accept arbitrary BJJ points as public keys \textbf{must} add an
  explicit subgroup check $[\ell]\pk_j = \mathcal{O}$ in
  \texttt{CreateDealing} before the DH computation.
\end{remark}

% ── 6.2 ──────────────────────────────────────────────────────────────────────
\subsection{Round 0: Non-Interactive Dealing}
\label{sec:dealing}

Each dealer $d \in [1, n]$ performs the following:

\begin{algorithm}[H]
\caption{$\mathsf{CreateDealing}(d, \{\pk_i\}_{i \ne d}, \sid, t, n)$}
\label{alg:dealing}
\begin{algorithmic}[1]
\State \textbf{Require:} $\mathsf{CurveSuite} \ne \bot$ \Comment{$\mathsf{ErrNilSuite}$}
\State \textbf{Validate:} $d \in [1, n]$; $t \ge 2$; $n \ge t$; $n \le 256$; all IDs unique in $[1, n]$

\Statex \hrulefill\ \textit{Polynomial \& VSS}\ \hrulefill
\State $\omega_d \getsr \Fr$ \Comment{Secret for this dealer}
\State $f_d(x) \leftarrow \omega_d + a_1 x + \cdots + a_{t-1} x^{t-1}$,
  \quad $a_i \getsr \Fr$
\State $A_k \leftarrow [a_k]\genbn$ for $k = 0, \ldots, t-1$
  \Comment{VSS commitments}
\State $s_i \leftarrow f_d(i)$ for $i = 1, \ldots, n$
  \Comment{Shamir shares over $\Fr$}

\Statex \hrulefill\ \textit{Session Data}\ \hrulefill
\State $\mathsf{randomMsg}_d \getsr \{0,1\}^{256}$
\State $\pi_{\mathrm{id}} \leftarrow \mathsf{ProveIdentity}(\sk_d, \pk_d, \sid)$
  \Comment{Schnorr PoK}
\State $\alpha \leftarrow \hash_r(\texttt{"golden-lhl-alpha"} \,\|\, \sid)$
  \Comment{LHL coefficient}

\Statex \hrulefill\ \textit{Share Encryption}\ \hrulefill
\For{each peer $j \ne d$}
  \State $(\pad_{d,j},\, R_{d,j}) \leftarrow
    \mathsf{DerivePad}(\sk_d, \pk_j, (\sid, \mathsf{randomMsg}_d), \alpha)$
  \State $z_{d,j} \leftarrow \pad_{d,j} + s_j \in \Fr$
    \Comment{Encrypt share}
  \State $\pi_{d,j} \leftarrow \PLONK.\Prove(\mathsf{EVRFCircuit},\, \sk_d, \ldots)$
    \Comment{ZK proof}
  \State $\mathsf{Erase}(\pad_{d,j})$
\EndFor

\Statex \hrulefill\ \textit{Cleanup}\ \hrulefill
\State $\mathsf{Erase}(f_d,\, \omega_d,\, \{s_j\}_{j \ne d})$

\State \Return $\mathsf{Round0Msg} = \bigl(\sid,\, d,\, \mathsf{randomMsg}_d,\,
  \mathbf{A},\, \{(R_{d,j}, z_{d,j})\}_{j \ne d},\, \pi_{\mathrm{id}},\,
  \{\pi_{d,j}\}_{j \ne d}\bigr)$
\end{algorithmic}
\end{algorithm}

\begin{remark}
  The dealer retains their own share $s_d = f_d(d)$ as private state for the
  completion phase. All other secret material (polynomial coefficients, peer
  shares, pads) is securely erased.
\end{remark}

% ── 6.3 ──────────────────────────────────────────────────────────────────────
\subsection{Dealing Verification}
\label{sec:verify}

Each verifier $v$ checks a dealing message from dealer $d$:

\begin{algorithm}[H]
\caption{$\mathsf{VerifyDealing}(\mathsf{msg},\, \pk_d,\, \{\pk_j\}_{j \ne d},\, \sid,\, t)$}
\label{alg:verify}
\begin{algorithmic}[1]
\State \textbf{Require:} $\mathsf{CurveSuite} \ne \bot$ \Comment{$\mathsf{ErrNilSuite}$}
\State \textbf{Check:} $\mathsf{msg.SessionID} = \sid$
\State \textbf{Check:} $|\mathbf{A}| = t$
  \Comment{VSS commitment count matches threshold}
\State \textbf{Check:} $A_0 \ne \mathcal{O}$
  \Comment{Non-degenerate PK contribution}
\State \textbf{Validate recipient set:}
  (a)~$\pi_{\mathrm{id}} \ne \bot$;
  (b)~dealer $d \notin \mathsf{Ciphertexts}.\mathsf{keys}$;
  (c)~$|\mathsf{Ciphertexts}| = |\mathsf{Proofs}| = n - 1$;
  (d)~$\forall\, j \ne d$: $j \in \mathsf{Ciphertexts}.\mathsf{keys} \cap \mathsf{Proofs}.\mathsf{keys}$;
  (e)~$\forall\, j$: $\mathsf{Ciphertexts}[j] \ne \bot$;
  (f)~$\forall\, j$: $R_{d,j} \ne \mathcal{O}$

\Statex \hrulefill\ \textit{Identity Proof}\ \hrulefill
\State $\mathsf{VerifyIdentity}(\pk_d, \sid, \pi_{\mathrm{id}})$
  \Comment{Schnorr PoK on BJJ}

\Statex \hrulefill\ \textit{Alpha Derivation}\ \hrulefill
\State $\alpha \leftarrow \hash_r(\texttt{"golden-lhl-alpha"} \,\|\, \sid)$

\Statex \hrulefill\ \textit{Ciphertext Consistency}\ \hrulefill
\For{each recipient $j$}
  \State Compute $\mathsf{LHS} = [z_{d,j}]\genbn$
  \State Compute $\mathsf{RHS} = R_{d,j} + \VSS(j)$
    \Comment{$\VSS(j) = \sum_k [j^k] A_k$}
  \State \textbf{Check:} $\mathsf{LHS} = \mathsf{RHS}$
    \Comment{See Proposition~\ref{prop:ct-verify}}
\EndFor

\Statex \hrulefill\ \textit{eVRF Proofs}\ \hrulefill
\For{each recipient $j$}
  \State Recompute $H_1, H_2$ from $\mathsf{sessData} = (\sid, \mathsf{randomMsg}_d)$
  \State $\PLONK.\Verify(\pi_{d,j},\, \mathsf{vk},\, (\pk_d, \pk_j, H_1, H_2, \alpha, R_{d,j}))$
\EndFor

\State \Return $\mathsf{accept}$
\end{algorithmic}
\end{algorithm}

\begin{remark}[Verifier Scope]
  Each verifier checks \emph{all} $n - 1$ ciphertext consistency equations and
  \emph{all} $n - 1$ eVRF proofs in a dealing, not just the one addressed to
  itself. This ensures a malicious dealer cannot selectively cheat by providing
  a correct ciphertext/proof for one recipient while providing an incorrect one
  for another.
\end{remark}

\begin{proposition}[Ciphertext Consistency]
  \label{prop:ct-verify}
  The verification equation $[z_{d,j}]\genbn = R_{d,j} + \VSS(j)$ holds for an
  honest dealer:
  \begin{align}
    [z_{d,j}]\genbn
    &= [\pad_{d,j} + s_j]\genbn \\
    &= [\pad_{d,j}]\genbn + [s_j]\genbn \\
    &= R_{d,j} + [f_d(j)]\genbn \\
    &= R_{d,j} + \VSS(j) \qquad \text{(by~\eqref{eq:vss-verify})}
  \end{align}
\end{proposition}

\medskip
\noindent\textbf{Worked example (toy field $\F_{13}$).}
Consider a group of order $13$ with generator $G$, threshold $t = 2$, and
a dealer with polynomial $f(x) = 5 + 3x$ (so $\omega = 5$, $a_1 = 3$).
The VSS commitments are $A_0 = [5]G$, $A_1 = [3]G$.
For recipient $j = 2$, the share is $s_2 = f(2) = 5 + 6 = 11 \bmod 13$.
Suppose the pad is $\pad = 7$. Then:
\begin{itemize}[nosep]
  \item $z_{d,2} = 7 + 11 = 18 \equiv 5 \pmod{13}$.
  \item $R_{d,2} = [7]G$.
  \item $\VSS(2) = [2^0]A_0 + [2^1]A_1 = [5]G + [6]G = [11]G$.
  \item $\mathsf{LHS} = [5]G$.
  \item $\mathsf{RHS} = [7]G + [11]G = [18]G = [5]G$. \checkmark
\end{itemize}

\begin{remark}
  The ciphertext consistency check is \emph{independent} of the eVRF proof.
  Together, they guarantee: (1)~the encrypted share is consistent with the
  committed polynomial (ciphertext consistency), and (2)~the pad commitment was
  correctly derived from the DH shared secret (eVRF proof). A dishonest dealer
  cannot produce a valid ciphertext with a wrong share, nor a valid proof with
  a wrong pad.
\end{remark}

% ── 6.4 ──────────────────────────────────────────────────────────────────────
\subsection{Protocol Completion}
\label{sec:complete}

After receiving and verifying all $n-1$ peer dealings, participant $i$
aggregates shares and computes the group key.

\begin{remark}[Pre-Verified Dealings]
  \label{rem:pre-verified}
  \texttt{Complete} assumes that every dealing in the input set has
  \emph{already passed} \texttt{VerifyDealing} (Algorithm~\ref{alg:verify}).
  It does \textbf{not} re-run verification internally. Passing an unverified or
  invalid dealing to \texttt{Complete} is a caller error that can produce
  incorrect key shares. The reference implementation's \texttt{Complete}
  function documents this precondition; callers are responsible for filtering
  out invalid dealings before aggregation.
\end{remark}

\begin{algorithm}[H]
\caption{$\mathsf{Complete}(i,\, \mathsf{ownDealing},\, \{\mathsf{msg}_d\}_{d \ne i},\,
  \{\pk_d\}_{d \ne i},\, \sid,\, t,\, n)$}
\label{alg:complete}
\begin{algorithmic}[1]
\State \textbf{Require:} $\mathsf{CurveSuite} \ne \bot$ \Comment{$\mathsf{ErrNilSuite}$}
\State \textbf{Require:} All peer dealings have passed $\mathsf{VerifyDealing}$
  \Comment{Caller precondition}
\State \textbf{Validate:}
  (a)~$|\mathsf{peerDealings}| = n - 1$;
  (b)~no dealing from self ($d \ne i$ for all dealers $d$);
  (c)~all dealer IDs unique;
  (d)~all dealer IDs from known participants;
  (e)~all session IDs match $\sid$

\Statex \hrulefill\ \textit{Share Aggregation}\ \hrulefill
\State $\mathsf{secretShare}_i \leftarrow s_{i,\mathrm{own}}$
  \Comment{Own share from own polynomial}
\State $\alpha \leftarrow \hash_r(\texttt{"golden-lhl-alpha"} \,\|\, \sid)$

\For{each peer dealing from dealer $d$}
  \State Retrieve ciphertext $(R_{d,i}, z_{d,i})$
  \State $(\pad_{d,i}, \_) \leftarrow
    \mathsf{DerivePad}(\sk_i, \pk_d, (\sid, \mathsf{randomMsg}_d), \alpha)$
    \Comment{Re-derive pad}
  \State $\hat{s}_{d,i} \leftarrow z_{d,i} - \pad_{d,i} \in \Fr$
    \Comment{Decrypt share}
  \State $\mathsf{Erase}(\pad_{d,i})$
  \State $\mathsf{secretShare}_i \leftarrow \mathsf{secretShare}_i + \hat{s}_{d,i}$
    \Comment{Accumulate in $\Fr$}
  \State $\mathsf{Erase}(\hat{s}_{d,i})$
\EndFor

\Statex \hrulefill\ \textit{Group Key}\ \hrulefill
\State $\mathsf{GroupKey} \leftarrow \sum_{d=1}^{n} A_{d,0}$
  \Comment{Sum of $[\omega_d]\genbn$ over all dealers}

\Statex \hrulefill\ \textit{Public Key Shares}\ \hrulefill
\For{$j = 1, \ldots, n$}
  \State $\mathsf{PKShare}_j \leftarrow \sum_{d=1}^{n} \VSS_d(j)$
    \Comment{$= [\mathsf{secretShare}_j]\genbn$ for honest execution}
\EndFor

\State \Return $(\mathsf{GroupKey},\, \{\mathsf{PKShare}_j\}_{j=1}^n,\, \mathsf{secretShare}_i)$
\end{algorithmic}
\end{algorithm}

\begin{theorem}[DKG Correctness]
  \label{thm:dkg-correct}
  After honest execution of the GOLDEN DKG:
  \begin{enumerate}[nosep]
    \item The combined polynomial is
      $F(x) = \sum_{d=1}^{n} f_d(x)$, of degree $t-1$ over $\Fr$.
    \item The group secret is
      $\Omega = F(0) = \sum_{d=1}^{n} \omega_d \in \Fr$.
    \item The group key is
      $\mathsf{GroupKey} = [\Omega]\genbn$.
    \item Each participant's aggregated share is
      $\mathsf{secretShare}_i = F(i) = \sum_{d=1}^{n} f_d(i) \in \Fr$.
    \item Any $t$ shares can reconstruct $\Omega$ via Lagrange
      interpolation~\eqref{eq:lagrange}.
    \item The public key share of participant $j$ satisfies
      $\mathsf{PKShare}_j = [\mathsf{secretShare}_j]\genbn$.
  \end{enumerate}
\end{theorem}

\begin{proof}
  \textit{(1)} follows from the definition: each $f_d$ is degree $t-1$, so
  their sum $F = \sum f_d$ is also degree $t-1$.

  \textit{(2)} $F(0) = \sum f_d(0) = \sum \omega_d$ by linearity.

  \textit{(3)} $\mathsf{GroupKey} = \sum A_{d,0} = \sum [\omega_d]\genbn
  = [\sum \omega_d]\genbn = [\Omega]\genbn$.

  \textit{(4)} By linearity of polynomial evaluation:
  $\mathsf{secretShare}_i = s_{i,\mathrm{own}} + \sum_{d \ne i} \hat{s}_{d,i}
  = f_i(i) + \sum_{d \ne i} f_d(i) = \sum_{d=1}^n f_d(i) = F(i)$.

  \textit{(5)} $F$ is degree $t-1$, so $t$ evaluations determine $F$
  uniquely by Lagrange interpolation.

  \textit{(6)} $\sum_d \VSS_d(j) = \sum_d [f_d(j)]\genbn
  = [\sum_d f_d(j)]\genbn = [F(j)]\genbn = [\mathsf{secretShare}_j]\genbn$.
\end{proof}

% ── 6.5 ──────────────────────────────────────────────────────────────────────
\subsection{Multi-Curve Derived Shares}
\label{sec:derived}

The GOLDEN DKG supports generating key shares for \emph{multiple} elliptic
curve groups from a single DKG execution. This allows a single $n$-party DKG
to simultaneously produce, e.g., BN254~G1 key shares (for the primary FROST
signing group) together with Baby~Jubjub or secp256k1 key shares for other
applications, without running separate DKG instances.

\subsubsection{Motivation}

Running independent DKGs for each target curve would require $k$ separate
broadcast rounds and $k \times n \times (n-1)$ PLONK proofs. By deriving
shares from the \emph{same} polynomial, the protocol produces $k$ sets of
correlated key shares with only $n \times (n-1)$ proofs (the primary eVRF
proofs). The derived curves rely on algebraic consistency checks rather than
additional zero-knowledge proofs.

\subsubsection{Coefficient Derivation}

Let $f_d(x) = a_0 + a_1 x + \cdots + a_{t-1} x^{t-1}$ be the dealer's
polynomial over $\Fr$ (the primary outer-group field). For each derived group
$\G_d$ with scalar field $\F_{q_d}$, the derived polynomial coefficients are:
\begin{equation}
  a'_k = a_k \bmod q_d, \qquad k = 0, \ldots, t-1
  \label{eq:derived-coeff}
\end{equation}
where $a_k$ is the canonical byte serialization of the $\Fr$ coefficient,
interpreted as an integer and reduced modulo $q_d$ (via \texttt{SetBytes},
which auto-reduces).

\begin{remark}[Statistical Bias]
  When $q_d < r$ (e.g., BJJ where $\ell \approx 2^{251}$ and $r \approx 2^{254}$),
  the reduction $a_k \bmod \ell$ introduces a statistical bias. A uniformly
  random $a_k \in [0, r)$ reduced modulo $\ell$ produces a distribution with
  statistical distance at most $\lfloor r / \ell \rfloor / r \approx 2^{-248}$
  from uniform over $\F_\ell$. This is cryptographically negligible and does
  not affect the Shamir threshold property. When $q_d > r$ (e.g., secp256k1
  where $q \approx 2^{256}$), no reduction occurs; coefficients are injected
  exactly.
\end{remark}

The derived polynomial is $f'_d(x) = a'_0 + a'_1 x + \cdots + a'_{t-1} x^{t-1}$
over $\F_{q_d}$. Shares $s'_j = f'_d(j)$ and VSS commitments
$A'_k = [a'_k] G_d$ are computed independently on the derived group.

\subsubsection{Derived Ciphertexts}

For each peer $j$, the cached pad bytes from the primary eVRF pad derivation
(the canonical byte representation of $\pad \in \Fr$) are reduced modulo
$q_d$:
\begin{equation}
  \pad'_{d,j} = \pad_{d,j} \bmod q_d
  \label{eq:derived-pad}
\end{equation}
The derived ciphertext is:
\begin{equation}
  z'_{d,j} = \pad'_{d,j} + s'_j \in \F_{q_d}
  \label{eq:derived-ct}
\end{equation}
An independent R~commitment is computed on the derived group:
\begin{equation}
  R'_{d,j} = [\pad'_{d,j}] G_d \in \G_d
  \label{eq:derived-rcommit}
\end{equation}
Note that the derived R~commitment is \emph{not} the reduction of the primary
$R_{d,j}$; it is freshly computed on $\G_d$ from the reduced pad. During
verification, if $R'_{d,j} = \mathcal{O}$ the dealing is rejected with
$\mathsf{ErrIdentityRCommitment}$ (same check as the primary R~commitment).

\subsubsection{Derived Verification}

Verification of derived ciphertexts uses the same algebraic consistency check
as the primary curve:
\begin{equation}
  [z'_{d,j}] G_d \stackrel{?}{=} R'_{d,j} + \VSS'_d(j)
  \label{eq:derived-verify}
\end{equation}
where $\VSS'_d(j) = \sum_{k=0}^{t-1} [j^k] A'_k$. The correctness argument
is identical to Proposition~\ref{prop:ct-verify} with all operations in $\G_d$.

\begin{proposition}[Derived Share Confidentiality]
  \label{prop:derived-conf}
  Let $\G_d$ be a derived group with scalar field $\F_{q_d}$ where $q_d \mid r$
  (the BN254 scalar field order). If the primary pad $\pad_{d,j} \in \Fr$ is
  computationally indistinguishable from uniform (Theorem~\ref{thm:pad-ind}),
  then the derived pad $\pad'_{d,j} = \pad_{d,j} \bmod q_d$ is computationally
  indistinguishable from uniform in $\F_{q_d}$.
\end{proposition}

\begin{proof}
  Since $\pad_{d,j}$ is computationally indistinguishable from $U \getsr \Fr$,
  it suffices to show that $U \bmod q_d$ is statistically close to uniform in
  $\F_{q_d}$. Write $r = k \cdot q_d + \rho$ with $0 \le \rho < q_d$. The
  reduction $U \bmod q_d$ is uniform over the residues $0, \ldots, q_d - 1$
  with a statistical bias of at most $\rho / r \le q_d / r < 2^{-\lambda}$
  (for $\lambda \approx 100$ when $q_d \approx 2^{253}$ and $r \approx 2^{254}$).
  The computational indistinguishability of $\pad_{d,j}$ from $U$ adds at most
  $\negl(\lambda)$ advantage, so the derived pad is computationally
  indistinguishable from uniform in $\F_{q_d}$.
\end{proof}

\begin{remark}[Security of Derived Share Encryption]
  No additional eVRF proofs are generated for derived curves. The security of
  derived share encryption follows from Proposition~\ref{prop:derived-conf}
  (derived pad indistinguishability) combined with the algebraic consistency
  check~\eqref{eq:derived-verify} (which ties the ciphertext to the derived
  VSS polynomial). A malicious dealer who passes the primary eVRF check and
  all derived consistency checks has committed to polynomials whose evaluations
  are correctly encrypted under the derived pads.
\end{remark}

\subsubsection{Derived Completion}

During completion (Algorithm~\ref{alg:complete}), each participant aggregates
derived shares analogously to the primary shares: start with the own derived
share, then decrypt each peer's derived ciphertext by re-deriving the pad bytes
(cached from the primary $\mathsf{DerivePad}$ call), reducing modulo $q_d$,
and subtracting. The derived group key, public key shares, and secret share are
computed identically to the primary output, but over $\G_d$.

The derived output is converted to a FROST key share via
$\mathsf{DkgOutputToDerivedKeyShare}$ (Section~\ref{sec:derived-bridge}).

% ══════════════════════════════════════════════════════════════════════════════
\section{FROST Integration}
\label{sec:frost}

The GOLDEN DKG output is designed for direct compatibility with the FROST
threshold Schnorr signature scheme (RFC~9591).

% ── 7.1 ──────────────────────────────────────────────────────────────────────
\subsection{Key Share Conversion}
\label{sec:bridge}

The bridge function $\mathsf{DkgOutputToKeyShare}$ maps:

\begin{center}
\renewcommand{\arraystretch}{1.3}
\begin{tabular}{@{}ll@{}}
\toprule
\textbf{FROST Field} & \textbf{GOLDEN DKG Source} \\
\midrule
$\mathsf{KeyShare.ID}$ & $i$ encoded as big-endian \texttt{uint32} at bytes $[28{:}32]$ \\
$\mathsf{KeyShare.SecretKey}$ & $\mathsf{secretShare}_i \in \Fr$ \\
$\mathsf{KeyShare.PublicKey}$ & $\mathsf{PKShare}_i \in \Gbn$ \\
$\mathsf{KeyShare.GroupKey}$ & $\mathsf{GroupKey} \in \Gbn$ \\
\bottomrule
\end{tabular}
\end{center}

The ID encoding matches FROST's internal \texttt{scalarFromInt} function:
a 32-byte zero-padded buffer with the node ID as a big-endian \texttt{uint32}
in the last 4 bytes. This ensures Lagrange coefficient computation in FROST
is consistent with the Shamir evaluation points used in the DKG.

\subsubsection{Derived Key Share Conversion}
\label{sec:derived-bridge}

For each derived group $\G_d$, the function
$\mathsf{DkgOutputToDerivedKeyShare}(\G_d, i, \mathsf{DerivedOutput})$
performs the same conversion as $\mathsf{DkgOutputToKeyShare}$ but using the
derived group's generator and scalar field. The ID encoding is identical
(big-endian \texttt{uint32} at bytes $[28{:}32]$). The resulting FROST key
share can be used for threshold signing on the derived group.

% ── 7.2 ──────────────────────────────────────────────────────────────────────
\subsection{FROST Threshold Schnorr Signing}
\label{sec:frostsign}

Given $t$ signers with key shares from the DKG, FROST produces a standard
Schnorr signature verifiable against the group key. The protocol has two
rounds:

\subsubsection{Round 1: Commitment}

Each signer $i$ in the signing set $\mathcal{S}$ ($|\mathcal{S}| \ge t$):
\begin{enumerate}[nosep]
  \item Sample $d_i, e_i \getsr \Fr \setminus \{0\}$ \hfill (hiding and binding nonces)
  \item $D_i \leftarrow [d_i]\genbn$, \quad $E_i \leftarrow [e_i]\genbn$
    \hfill (nonce commitments)
  \item Broadcast $(\id_i, D_i, E_i)$
\end{enumerate}

\subsubsection{Round 2: Signature Share}

\begin{remark}[Notation for $z$]
  The symbol $z$ is used for both DKG encrypted shares ($z_{d,j}$ in
  Section~\ref{sec:encrypt}) and FROST signature shares ($z_i$ below). The
  subscript convention distinguishes them: $z_{d,j}$ is indexed by
  dealer--recipient pair, while $z_i$ is indexed by signer. In the aggregated
  signature, $z$ (no subscript) denotes the combined signature scalar.
\end{remark}

Each signer $i$, given all commitments $\{(\id_j, D_j, E_j)\}_{j \in \mathcal{S}}$
and message $m$:
\begin{enumerate}[nosep]
  \item Encode commitments:
    $\mathsf{enc} = |\mathcal{S}| \,\|\, (\id_1 \,\|\, D_1 \,\|\, E_1) \,\|\, \cdots$
    (length-prefixed)
  \item For each signer $j$: $\rho_j = \hash_r^{(1)}(m, \mathsf{enc}, \id_j)$
    \hfill (binding factor)
  \item Group commitment: $R = \sum_{j \in \mathcal{S}} (D_j + [\rho_j] E_j)$
  \item Challenge: $c = \hash_r^{(2)}(R, \mathsf{GroupKey}, m)$
  \item Lagrange coefficient:
    $\lambda_i = \prod_{j \in \mathcal{S},\, j \ne i} \frac{\id_j}{\id_j - \id_i}$
  \item Signature share: $z_i = d_i + \rho_i \cdot e_i + \lambda_i \cdot \mathsf{secretKey}_i \cdot c$
  \item Securely erase $d_i, e_i$
\end{enumerate}

\subsubsection{Aggregation}

The coordinator (or any party) computes:
\begin{equation}
  z = \sum_{i \in \mathcal{S}} z_i, \qquad \sigma = (R, z)
  \label{eq:aggregate}
\end{equation}

\subsubsection{Verification}

Standard Schnorr verification:
\begin{equation}
  [z]\genbn \stackrel{?}{=} R + [c]\mathsf{GroupKey}
  \label{eq:verify-sig}
\end{equation}
where $c = \hash_r^{(2)}(R, \mathsf{GroupKey}, m)$.

\begin{proposition}[FROST Correctness]
  For an honestly generated signature:
  \begin{align}
    [z]\genbn
    &= \Bigl[\sum_{i} \bigl(d_i + \rho_i e_i + \lambda_i \cdot \mathsf{sk}_i \cdot c\bigr)\Bigr]\genbn \\
    &= \sum_{i} \bigl([d_i]\genbn + [\rho_i][e_i]\genbn\bigr)
      + c \cdot \Bigl[\sum_{i} \lambda_i \cdot \mathsf{sk}_i\Bigr]\genbn \\
    &= \sum_{i} (D_i + [\rho_i]E_i) + c \cdot [F(0)]\genbn \\
    &= R + [c]\mathsf{GroupKey}
  \end{align}
  The third equality uses the Lagrange interpolation property:
  $\sum_{i \in \mathcal{S}} \lambda_i \cdot F(i) = F(0) = \Omega$.
\end{proposition}

% ── 7.3 ──────────────────────────────────────────────────────────────────────
\subsection{Piggyback Nonce Preprocessing}
\label{sec:piggyback}

Standard FROST signing (Section~\ref{sec:frostsign}) requires two interactive
rounds: Round~1 broadcasts nonce commitments, and Round~2 broadcasts signature
shares.  Piggyback nonce preprocessing reduces the steady-state online cost to
\emph{one round} by attaching each signer's \emph{next} nonce commitment
alongside the current signature share.  A one-time two-round bootstrap
initializes the nonce pipeline; every subsequent session is a single message.

% ── 7.3.1 ────────────────────────────────────────────────────────────────────
\subsubsection{Motivation}
\label{sec:piggyback-motivation}

In latency-sensitive applications (e.g., blockchain validation, real-time
co-signing), eliminating a round trip is significant.  The core observation is
that FROST Round~1 generates nonces \emph{independently} of the message to be
signed---only Round~2 depends on the message and on other signers'
commitments.  Therefore a signer can pre-generate its Round~1 nonce for
session~$N{+}1$ during session~$N$ and broadcast the commitment alongside the
Round~2 share for session~$N$.  Recipients store the piggybacked commitment
and use it as the Round~1 commitment in session~$N{+}1$, skipping that round
entirely.

After a two-round bootstrap that initializes the nonce pipeline, every
subsequent session exchanges a single message per signer:
$(z_i,\, D_i^{\mathrm{next}},\, E_i^{\mathrm{next}})$.

% ── 7.3.2 ────────────────────────────────────────────────────────────────────
\subsubsection{Piggyback State}
\label{sec:piggyback-state}

Each signer $i$ maintains a piggyback state tuple:
\begin{equation}
  \mathsf{PBState}_i = (\id_i,\; d^{\mathrm{next}},\; e^{\mathrm{next}},\;
    D^{\mathrm{next}},\; E^{\mathrm{next}},\; \sigma,\; \mathsf{active})
  \label{eq:pbstate}
\end{equation}
where:
\begin{itemize}[nosep]
  \item $\id_i$: the signer's participant identifier.
  \item $d^{\mathrm{next}}, e^{\mathrm{next}} \in \Fr \setminus \{0\}$:
    pre-generated hiding and binding nonce scalars for the next session
    (secret; zeroed after consumption).
  \item $D^{\mathrm{next}} = [d^{\mathrm{next}}]\genbn$,\;
    $E^{\mathrm{next}} = [e^{\mathrm{next}}]\genbn$:
    corresponding public commitments (broadcast with the current share).
  \item $\sigma \in \{0, 1, \ldots, 2^{64}-1\}$: a monotonically increasing
    session index, cryptographically bound into the binding factor to prevent
    cross-session commitment replay.
  \item $\mathsf{active} \in \{\mathsf{true}, \mathsf{false}\}$: whether the
    state holds a valid pending nonce ready for consumption.
\end{itemize}

\begin{definition}[Session-Bound Message]
  \label{def:session-msg}
  For session index $\sigma$ and application message $m$, the session-bound
  message is:
  \begin{equation}
    \mathsf{sessionMsg} = \mathrm{be64}(\sigma) \,\|\, m
    \label{eq:session-msg}
  \end{equation}
  where $\mathrm{be64}$ encodes $\sigma$ as 8 big-endian bytes.  This
  substitutes for $m$ in all $\mathsf{SignRound2}$ calls within the
  piggyback protocol, ensuring the binding factor
  $\rho_j = \hash_r^{(1)}(\mathsf{sessionMsg}, \mathsf{enc}, \id_j)$ is
  session-specific.
\end{definition}

% ── 7.3.3 ────────────────────────────────────────────────────────────────────
\subsubsection{Bootstrap Protocol}
\label{sec:piggyback-bootstrap}

The bootstrap is a standard two-round FROST session augmented with
pre-generation of the next nonce pair.  It transitions the state from
$\mathsf{INACTIVE}$ to $\mathsf{ACTIVE}$.

\begin{algorithm}[H]
\caption{$\mathsf{BootstrapRound1}(r,\, \mathsf{state},\, \mathsf{share})$}
\label{alg:bootstrap-r1}
\begin{algorithmic}[1]
\State \textbf{Lock} $\mathsf{state}.\mu$
\Statex \hrulefill\ \textit{Session 0 nonce}\ \hrulefill
\State $(d_0, e_0, D_0, E_0) \leftarrow \mathsf{SignRound1}(r, \mathsf{share})$
  \Comment{Current session nonce}
\Statex \hrulefill\ \textit{Session 1 nonce (pre-generated)}\ \hrulefill
\State $(d^{\mathrm{next}}, e^{\mathrm{next}}, D^{\mathrm{next}}, E^{\mathrm{next}})
  \leftarrow \mathsf{SignRound1}(r, \mathsf{share})$
\If{RNG failure}
  \State $d_0.\mathsf{Zero}()$;\; $e_0.\mathsf{Zero}()$
  \State \Return $\bot$
\EndIf
\State $\mathsf{state}.(d^{\mathrm{next}}, e^{\mathrm{next}}, D^{\mathrm{next}}, E^{\mathrm{next}})
  \leftarrow (d^{\mathrm{next}}, e^{\mathrm{next}}, D^{\mathrm{next}}, E^{\mathrm{next}})$
\State $\mathsf{state}.\sigma \leftarrow 1$
\Comment{Next session is $\sigma = 1$}
\State \textbf{Unlock} $\mathsf{state}.\mu$
\State \Return $(d_0, e_0, D_0, E_0)$
  \Comment{Caller broadcasts $(D_0, E_0)$}
\end{algorithmic}
\end{algorithm}

After collecting all signers' commitments for session~0, each signer runs:

\begin{algorithm}[H]
\caption{$\mathsf{BootstrapRound2}(\mathsf{share},\, \mathsf{state},\,
  (d_0, e_0),\, m,\, \{(\id_j, D_j, E_j)\}_{j \in \mathcal{S}})$}
\label{alg:bootstrap-r2}
\begin{algorithmic}[1]
\State \textbf{Lock} $\mathsf{state}.\mu$
\State $\mathsf{sessionMsg} \leftarrow \mathrm{be64}(0) \,\|\, m$
  \Comment{Bootstrap is session $\sigma = 0$}
\State $z_i \leftarrow \mathsf{SignRound2}(\mathsf{share},\, (d_0, e_0),\,
  \mathsf{sessionMsg},\, \mathsf{commitments})$
\Comment{$d_0, e_0$ zeroed by \texttt{defer}}
\State $\mathsf{state}.\mathsf{active} \leftarrow \mathsf{true}$
\State \textbf{Unlock} $\mathsf{state}.\mu$
\State \Return $\bigl(z_i,\; \mathsf{state}.D^{\mathrm{next}},\;
  \mathsf{state}.E^{\mathrm{next}},\; \mathsf{state}.\sigma\bigr)$
  \Comment{Piggybacked share}
\end{algorithmic}
\end{algorithm}

% ── 7.3.4 ────────────────────────────────────────────────────────────────────
\subsubsection{Steady-State Protocol}
\label{sec:piggyback-steady}

After bootstrap, each subsequent session is a single-round protocol.  The
signer consumes the pending nonce, signs, and simultaneously pre-generates and
broadcasts the next nonce.

\begin{algorithm}[H]
\caption{$\mathsf{PiggybackSign}(r,\, \mathsf{share},\, \mathsf{state},\,
  m,\, \{(\id_j, D_j, E_j)\}_{j \in \mathcal{S}})$}
\label{alg:piggyback-sign}
\begin{algorithmic}[1]
\State \textbf{Lock} $\mathsf{state}.\mu$

\Statex \hrulefill\ \textit{State validation}\ \hrulefill
\State \textbf{Require} $\mathsf{state}.\mathsf{active} = \mathsf{true}$
  \hfill \Comment{Else $\mathsf{ErrPiggybackInactive}$}
\State \textbf{Require} $\mathsf{state}.\sigma < 2^{64} - 1$
  \hfill \Comment{Else $\mathsf{ErrSessionOverflow}$}
\State \textbf{Require} $\mathsf{state}.d^{\mathrm{next}} \ne \bot$

\Statex \hrulefill\ \textit{Pre-validation (retryable on failure)}\ \hrulefill
\State \textbf{Check} $|\mathcal{S}| \ge t$
  \Comment{Threshold}
\State \textbf{Check} $\forall\, j$: $D_j \ne \mathcal{O}$ and $E_j \ne \mathcal{O}$
  \Comment{Identity point rejection (RFC~9591~\S4.3)}
\State \textbf{Check} no duplicate $\id_j$ in $\mathcal{S}$
\State \textbf{Check} $(\id_i, \mathsf{state}.D^{\mathrm{next}}, \mathsf{state}.E^{\mathrm{next}})
  \in \{(\id_j, D_j, E_j)\}_{j \in \mathcal{S}}$
  \Comment{Own commitment present}
\If{any check fails} \Return error \Comment{State unchanged; caller may retry}
\EndIf

\Statex \hrulefill\ \textit{Pre-generate next nonce (retryable on failure)}\ \hrulefill
\State $(d^{\mathrm{new}}, e^{\mathrm{new}}, D^{\mathrm{new}}, E^{\mathrm{new}})
  \leftarrow \mathsf{SignRound1}(r, \mathsf{share})$
\If{RNG failure} \Return error \Comment{State unchanged; caller may retry}
\EndIf

\Statex \hrulefill\ \textit{Point of no return: consume nonce}\ \hrulefill
\State $(d^{\mathrm{clone}}, e^{\mathrm{clone}}) \leftarrow
  \mathsf{Clone}(\mathsf{state}.d^{\mathrm{next}},\,
  \mathsf{state}.e^{\mathrm{next}})$
  \Comment{Deep copy via $\mathsf{Scalar.Set}$}
\State $\mathsf{state}.d^{\mathrm{next}}.\mathsf{Zero}()$;\;
  $\mathsf{state}.e^{\mathrm{next}}.\mathsf{Zero}()$
  \Comment{Erase original}
\State $\mathsf{state}.(d^{\mathrm{next}}, e^{\mathrm{next}},
  D^{\mathrm{next}}, E^{\mathrm{next}}) \leftarrow
  (d^{\mathrm{new}}, e^{\mathrm{new}}, D^{\mathrm{new}}, E^{\mathrm{new}})$
  \Comment{Install new}

\Statex \hrulefill\ \textit{Sign with session-bound message}\ \hrulefill
\State $\mathsf{sessionMsg} \leftarrow \mathrm{be64}(\mathsf{state}.\sigma) \,\|\, m$
\State $z_i \leftarrow \mathsf{SignRound2}(\mathsf{share},\,
  (d^{\mathrm{clone}}, e^{\mathrm{clone}}),\,
  \mathsf{sessionMsg},\, \mathsf{commitments})$
  \Comment{$d^{\mathrm{clone}}, e^{\mathrm{clone}}$ zeroed by \texttt{defer}}
\If{$\mathsf{SignRound2}$ fails}
  \State $\mathsf{state}.d^{\mathrm{next}}.\mathsf{Zero}()$;\;
    $\mathsf{state}.e^{\mathrm{next}}.\mathsf{Zero}()$
    \Comment{Erase newly installed nonce}
  \State $\mathsf{state}.\mathsf{active} \leftarrow \mathsf{false}$
    \Comment{$\mathsf{FAILED}$ state: must $\mathsf{ReBootstrap}$}
  \State \Return $\bot$
\EndIf

\Statex \hrulefill\ \textit{Success}\ \hrulefill
\State $\mathsf{state}.\sigma \leftarrow \mathsf{state}.\sigma + 1$
\State \textbf{Unlock} $\mathsf{state}.\mu$
\State \Return $\bigl(z_i,\; \mathsf{state}.D^{\mathrm{next}},\;
  \mathsf{state}.E^{\mathrm{next}},\; \mathsf{state}.\sigma\bigr)$
\end{algorithmic}
\end{algorithm}

The state machine is:
\[
  \mathsf{INACTIVE}
  \xrightarrow{\mathsf{Bootstrap}}
  \mathsf{ACTIVE}
  \xrightarrow[\text{success}]{\mathsf{PiggybackSign}}
  \mathsf{ACTIVE}
  \xrightarrow[\mathsf{SignRound2}\text{ fail}]{\mathsf{PiggybackSign}}
  \mathsf{FAILED}
  \xrightarrow{\mathsf{ReBootstrap}}
  \mathsf{INACTIVE}
\]

% ── 7.3.5 ────────────────────────────────────────────────────────────────────
\subsubsection{Commitment Collection and Aggregation}
\label{sec:piggyback-collect}

After receiving piggybacked shares from all signers in session~$N$, a
collector (any party) extracts the commitments for session~$N{+}1$ and
aggregates the current session's signature.

\medskip\noindent
\textbf{Commitment collection.}
For each received piggybacked share
$(z_j,\, D_j^{\mathrm{next}},\, E_j^{\mathrm{next}},\, \sigma_j)$:
\begin{enumerate}[nosep]
  \item \textbf{Non-nil:} $D_j^{\mathrm{next}} \ne \bot$ and
    $E_j^{\mathrm{next}} \ne \bot$.
  \item \textbf{ID match:} the commitment's signer ID equals the signature
    share's signer ID.
  \item \textbf{Session index:} $\sigma_j$ equals the expected next session
    index $\sigma_{\mathrm{expected}}$.
  \item \textbf{Identity check:} $D_j^{\mathrm{next}} \ne \mathcal{O}$ and
    $E_j^{\mathrm{next}} \ne \mathcal{O}$ (RFC~9591~\S4.3).
  \item \textbf{No duplicates:} each signer ID appears at most once.
\end{enumerate}
Collected commitments are stored in a map keyed by signer ID bytes and
returned sorted by ID for deterministic binding factor computation.

\medskip\noindent
\textbf{Aggregation.}
$\mathsf{PiggybackAggregate}$ combines signature aggregation and commitment
collection:
\begin{enumerate}[nosep]
  \item Extract the standard $\mathsf{SignatureShare}$ values from the
    piggybacked shares.
  \item Compute $\mathsf{sessionMsg} = \mathrm{be64}(\sigma) \,\|\, m$.
  \item Aggregate using
    $\mathsf{AggregateWithVerification}(\mathsf{sessionMsg},\,
    \mathsf{commitments},\, \mathsf{sigShares},\, \mathsf{publicKeys},\,
    \mathsf{GroupKey})$.
    \textbf{Per-share verification is mandatory}
    (Section~\ref{sec:frostsign}); plain $\mathsf{Aggregate}$ must not be used,
    as it cannot detect commitment substitution by a man-in-the-middle.
  \item Collect next-session commitments via $\mathsf{CollectPiggybackCommitments}$.
  \item Return $(\sigma_{\mathrm{sig}},\, \mathsf{nextCommitments})$.
\end{enumerate}

% ── 7.3.6 ────────────────────────────────────────────────────────────────────
\subsubsection{Re-Bootstrap Conditions}
\label{sec:piggyback-rebootstrap}

The nonce pipeline must be restarted (all signers in the active set
re-bootstrap) when any of the following occurs:
\begin{enumerate}[nosep]
  \item \textbf{Signer set change.} A signer goes offline or a new signer
    joins.  Different commitment sets produce different binding factors, so
    cached commitments are incompatible.
  \item \textbf{Stale session index.} A signer receives commitments with a
    session index that does not match its own $\sigma$.
  \item \textbf{$\mathsf{SignRound2}$ failure.} After nonce consumption, the
    signing step fails.  The pending nonce is zeroed and the state transitions
    to $\mathsf{FAILED}$; the chain is irrecoverably broken.
  \item \textbf{Explicit call.} The application explicitly invokes
    $\mathsf{ReBootstrap}()$, which securely erases all pending nonce material
    and resets $\sigma$ to~0.
\end{enumerate}

\begin{remark}
  Re-bootstrap by \emph{any} signer in the active set requires \emph{all}
  signers to re-bootstrap.  The detection mechanism is: a signer that needs
  re-bootstrap sends a standard Round~1 commitment (without piggyback
  data).  Recipients detect this as a bootstrap message and re-enter
  bootstrap mode themselves.
\end{remark}

% ── 7.3.7 ────────────────────────────────────────────────────────────────────
\subsubsection{Security Argument}
\label{sec:piggyback-security}

\begin{definition}[EUF-CMA Security]
  \label{def:euf-cma}
  A threshold signature scheme is \emph{existentially unforgeable under
  chosen-message attack} (EUF-CMA) if no probabilistic polynomial-time
  adversary, controlling up to $t{-}1$ corrupted signers, can produce a valid
  signature on a message $m^*$ that was never signed by the honest signers,
  even after adaptively requesting signatures on messages of its choice.
\end{definition}

\begin{proposition}[Piggyback Security Reduction]
  \label{prop:piggyback-security}
  If FROST is EUF-CMA secure (Definition~\ref{def:euf-cma}) against a static
  $t{-}1$ adversary, then Piggyback FROST is EUF-CMA secure under the same
  assumptions plus authenticated message delivery of piggybacked shares.
\end{proposition}

\begin{proof}[Proof sketch]
  We reduce Piggyback FROST security to standard FROST security by showing
  that each piggyback session~$N$ is indistinguishable from a standard
  two-round FROST session in which Round~1 was delivered during session~$N{-}1$.

  \begin{enumerate}
    \item \textbf{Reduction to standard FROST.}
      In session~$N$, each signer uses the nonce $(d_i, e_i)$ that was
      pre-generated in session~$N{-}1$ and whose commitment $(D_i, E_i)$ was
      broadcast alongside session~$N{-}1$'s signature share.  The signing
      computation in session~$N$ is \emph{identical} to
      $\mathsf{SignRound2}(\mathsf{share},\, (d_i, e_i),\,
      \mathsf{sessionMsg},\, \mathsf{commitments})$.  No structural change is
      made to the FROST signing algorithm; only the \emph{timing} of Round~1
      delivery differs.

    \item \textbf{Nonce independence.}
      Each nonce pair $(d_i, e_i)$ is sampled from fresh randomness via
      $\mathsf{SignRound1}$.  There is no cryptographic chain or derivation
      linking consecutive nonces: nonce~$N{+}1$ is independent of nonce~$N$.
      Compromising one session's nonce reveals nothing about adjacent sessions.

    \item \textbf{Binding factor separation.}
      The piggybacked commitments $(D_j^{\mathrm{next}}, E_j^{\mathrm{next}})$
      for session~$N{+}1$ are \emph{not} included in session~$N$'s encoded
      commitment list.  They are carried in the message envelope but never
      passed to $\mathsf{encodeCommitments}$ or
      $\mathsf{computeBindingFactors}$.  Therefore they do not affect
      session~$N$'s binding factors $\rho_j$ or challenge $c$.

    \item \textbf{Session isolation.}
      The session index $\sigma$ is prepended to the message before binding
      factor computation (Definition~\ref{def:session-msg}).  Even if an
      adversary replays a commitment from session~$N$ into session~$N'$, the
      binding factor $\rho_j = \hash_r^{(1)}(\mathrm{be64}(\sigma') \,\|\, m,\,
      \mathsf{enc},\, \id_j)$ will differ (since $\sigma \ne \sigma'$),
      producing a distinct effective nonce for the same $(d_j, e_j)$.  This
      prevents cross-session nonce reuse.
  \end{enumerate}

  Combining these four observations, any forgery against Piggyback FROST
  implies a forgery against standard FROST in some session~$N$ (with Round~1
  delivered earlier), contradicting the assumed EUF-CMA security of FROST.
\end{proof}

\begin{remark}[Transport Authentication Requirement]
  \label{rem:transport-auth}
  The security reduction assumes \emph{authenticated delivery} of piggybacked
  shares.  Concretely, the transport layer (TLS, Noise protocol, etc.) must
  guarantee that a received $(z_j,\, D_j^{\mathrm{next}},\,
  E_j^{\mathrm{next}},\, \sigma_j)$ was indeed sent by signer~$j$.
  Without authentication, an active adversary could substitute piggybacked
  commitments, causing honest signers to use adversary-chosen nonce points in
  the next session.  Mandatory use of $\mathsf{AggregateWithVerification}$
  detects such substitution \emph{after the fact} (invalid shares will fail
  verification), but authenticated delivery prevents the attack entirely.
\end{remark}

\begin{remark}[Commitment Binding Separation]
  \label{rem:binding-separation}
  It is critical that the piggybacked next-session commitments
  $(D_j^{\mathrm{next}}, E_j^{\mathrm{next}})$ are \emph{inert} with respect
  to the current session's cryptographic computation.  They are never hashed
  into the current session's binding factors or challenge.  The separation
  ensures that a malicious signer cannot influence the current session's
  security by choosing adversarial next-session commitments.
\end{remark}

\begin{remark}[Adaptive Message Choice]
  \label{rem:adaptive-msg}
  In the piggyback model, a signer commits to nonces for session~$N{+}1$
  \emph{before} the message $m_{N+1}$ is known. An adversary who sees
  $(D_j^{\mathrm{next}}, E_j^{\mathrm{next}})$ could attempt to choose
  $m_{N+1}$ adaptively based on these commitments. This is precisely the
  setting covered by FROST's EUF-CMA security proof: FROST's binding factor
  mechanism (which hashes the message, all commitments, and participant IDs)
  ensures that the adversary cannot exploit knowledge of commitments to forge
  a signature. The reduction in Proposition~\ref{prop:piggyback-security}
  therefore holds even under adaptive message selection---the adversary's
  power in the piggyback setting is a strict subset of its power in the
  standard FROST CMA game.
\end{remark}

% ══════════════════════════════════════════════════════════════════════════════
\section{Security Analysis}
\label{sec:security}

% ── 8.1 ──────────────────────────────────────────────────────────────────────
\subsection{Security Model and Assumptions}
\label{sec:secmodel}

The security model is defined in Section~\ref{sec:secmodel-intro}. We
summarize the key points:

\begin{itemize}
  \item \textbf{Authenticated broadcast channel.} The protocol assumes all
    Round0Msg broadcasts are delivered authentically with eventual delivery.
    The Schnorr identity proof binds the dealing to the dealer's BJJ key, but
    message integrity is assumed from the transport layer.

  \item \textbf{Static corruption (no adaptive).} The adversary chooses which
    parties to corrupt before the protocol begins. Corrupted parties may deviate
    arbitrarily (Byzantine faults, rushing adversary). Adaptive corruption is
    \emph{not} considered.

  \item \textbf{Threshold adversary.} At most $t - 1$ participants are
    corrupted.

  \item \textbf{Computational assumptions.} DDH on $\Gbj$, DL on $\Gbn$,
    PLONK knowledge soundness (AGM + ROM).
\end{itemize}

% ── 8.2 ──────────────────────────────────────────────────────────────────────
\subsection{Security Properties}
\label{sec:secprops}

We state the four main security theorems with game-based reductions.
In each game, $\mathcal{A}$ denotes a PPT adversary corrupting at most
$t - 1$ parties (chosen statically before Round~0).

\begin{theorem}[Pad Indistinguishability]
  \label{thm:pad-ind}
  Under the DDH assumption on $\Gbj$, no PPT adversary $\mathcal{A}$ can
  distinguish the pad $\pad_{d,j}$ from a uniformly random element of $\Fr$
  with non-negligible advantage, for any honest dealer $d$ and honest
  recipient $j$.
\end{theorem}

\begin{proof}
  We construct a sequence of games.

  \smallskip\noindent\textbf{Game~0 (Real).}
  $\mathcal{A}$ receives the transcript of an honest execution, including
  all dealings from corrupted parties and the public components of honest
  dealings. The pad is $\pad_{d,j} = \Extract([\dhx]H_1) + \alpha \cdot
  \Extract([\dhx]H_2)$ where $\dhx = \Extract([\sk_d]\pk_j)$.

  \smallskip\noindent\textbf{Game~1 (DDH Switch).}
  Replace the DH shared secret $S = [\sk_d \sk_j]\genbjj$ with a random
  subgroup point $S^* \getsr \Gbj$. By DDH on $\Gbj$, $|\Pr[\mathcal{A}
  \text{ wins in Game~0}] - \Pr[\mathcal{A} \text{ wins in Game~1}]| \le
  \epsilon_{\mathrm{DDH}}$.

  \smallskip\noindent\textbf{Game~2 (LHL Extraction).}
  In Game~1, $\dhx^* = \Extract(S^*)$ is the x-coordinate of a uniformly random
  subgroup point, yielding HILL min-entropy $\ge \log\ell - 1 \ge 252$ bits.
  The points $P_1^* = [\dhx^*]H_1$ and $P_2^* = [\dhx^*]H_2$ are independent
  (since $H_1, H_2$ are outputs of independent hash-to-curve evaluations).
  By the computational Leftover Hash Lemma (Lemma~\ref{lem:clhl}), replacing
  $\pad$ with $U \getsr \Fr$ introduces at most $\epsilon_{\mathrm{LHL}}$
  distinguishing advantage.

  \smallskip\noindent
  Total advantage: $\mathrm{Adv}(\mathcal{A}) \le \epsilon_{\mathrm{DDH}}
  + \epsilon_{\mathrm{LHL}} = \negl(\lambda)$.
\end{proof}

\begin{theorem}[Share Confidentiality]
  \label{thm:share-conf}
  Under the DDH assumption on $\Gbj$ and PLONK knowledge soundness, no PPT
  adversary corrupting at most $t - 1$ parties can distinguish the encrypted
  share $z_{d,j}$ of an honest dealer $d$ to an honest recipient $j$ from a
  uniformly random $\Fr$ element.
\end{theorem}

\begin{proof}
  Consider the following game:

  \smallskip\noindent\textbf{Game~0 (Real).}
  $\mathcal{A}$ sees $z_{d,j} = \pad_{d,j} + s_j$ where $s_j = f_d(j)$.

  \smallskip\noindent\textbf{Game~1 (Random Pad).}
  Replace $\pad_{d,j}$ with $U \getsr \Fr$. By Theorem~\ref{thm:pad-ind},
  $|\Pr[\text{Game~0}] - \Pr[\text{Game~1}]| \le \negl(\lambda)$.

  \smallskip\noindent
  In Game~1, $z_{d,j} = U + s_j$ is a one-time pad encryption: for any
  fixed $s_j$, the distribution of $z_{d,j}$ is uniform over $\Fr$. The
  PLONK knowledge soundness ensures the dealer actually uses the pad
  derived from the DH secret (otherwise the eVRF proof would not verify),
  preventing a malicious dealer from substituting a known pad value.
\end{proof}

\begin{theorem}[Dealing Integrity]
  \label{thm:dealing-integrity}
  Under the discrete logarithm assumption on $\Gbn$ and PLONK knowledge
  soundness, no PPT adversary (as a malicious dealer) can produce a dealing
  that passes $\mathsf{VerifyDealing}$ while distributing shares inconsistent
  with a degree-$(t{-}1)$ polynomial.
\end{theorem}

\begin{proof}
  Suppose a malicious dealer $d^*$ produces a dealing that passes verification.

  \smallskip\noindent\textbf{Step~1 (Ciphertext consistency).}
  For each recipient $j$, verification checks
  $[z_{d^*,j}]\genbn = R_{d^*,j} + \VSS(j)$ (Proposition~\ref{prop:ct-verify}).
  This implies $z_{d^*,j} = r_{d^*,j} + f_{d^*}(j)$ where $r_{d^*,j}$ is
  the discrete log of $R_{d^*,j}$ and $f_{d^*}(x)$ is the polynomial committed
  by the VSS commitments $\{A_k\}$. (If $z_{d^*,j} \ne r_{d^*,j} + f_{d^*}(j)$,
  the check fails by the DL assumption on $\Gbn$.)

  \smallskip\noindent\textbf{Step~2 (eVRF binding).}
  The PLONK proof $\pi_{d^*,j}$ attests that $R_{d^*,j} = [\pad]\genbn$
  where $\pad$ is the correctly derived LHL pad from the DH shared secret.
  By knowledge soundness, the dealer cannot produce a valid proof for an
  $R_{d^*,j}$ that is not the pad commitment, except with negligible
  probability.

  \smallskip\noindent\textbf{Step~3 (Share extraction).}
  Combining Steps~1--2: $z_{d^*,j} = \pad_{d^*,j} + f_{d^*}(j)$ for the
  committed polynomial $f_{d^*}$. The decrypted share is
  $\hat{s}_j = z_{d^*,j} - \pad_{d^*,j} = f_{d^*}(j)$, which lies on the
  degree-$(t{-}1)$ polynomial. Any deviation is detected.
\end{proof}

\begin{theorem}[Session Binding]
  \label{thm:session-binding}
  Identity proofs and eVRF pads are bound to the session ID $\sid$. No PPT
  adversary can replay a dealing from session $\sid'$ into session $\sid \ne
  \sid'$ such that verification accepts.
\end{theorem}

\begin{proof}
  The binding operates through two mechanisms:

  \smallskip\noindent\textbf{(a) Schnorr identity proof.}
  The challenge includes the session ID:
  $c = \hash_\ell(\texttt{"golden-pki-pok"} \,\|\, \sid \,\|\, \cdots)$.
  A proof generated for $\sid'$ yields challenge $c' = \hash_\ell(\cdots
  \sid' \cdots) \ne c$ (by collision resistance of $\hash_\ell$). The
  verification equation $[s]\genbjj = K + [c]\pk_d$ fails since $c' \ne c$.

  \smallskip\noindent\textbf{(b) eVRF pad domain separation.}
  The LHL coefficient $\alpha = \hash_r(\texttt{"golden-lhl-alpha"} \,\|\,
  \sid)$ and hash-to-curve inputs $\mathsf{sessData} = (\sid,
  \mathsf{randomMsg})$ both include $\sid$. A replayed dealing would produce
  different $H_1, H_2$ and $\alpha$ values, yielding a different pad
  $\pad'_{d,j} \ne \pad_{d,j}$. The eVRF proof (bound to the original
  $R_{d,j} = [\pad_{d,j}]\genbn$) will not verify against the recomputed
  public inputs.
\end{proof}

% ── 8.3 ──────────────────────────────────────────────────────────────────────
\subsection{Threshold Security}
\label{sec:threshold}

\begin{proposition}[Information-Theoretic Secrecy of Shares]
  Any $t - 1$ or fewer Shamir shares of a degree-$(t-1)$ polynomial over $\Fr$
  reveal no information about the constant term $\Omega = F(0)$. This is a
  standard property of Shamir secret sharing~\cite{shamir1979}.
\end{proposition}

\begin{proposition}[Threshold Signing]
  Any set of $t$ or more FROST signers can produce a valid Schnorr signature
  verifiable against $\mathsf{GroupKey} = [\Omega]\genbn$. Any $t - 1$ or fewer
  signers cannot produce a valid signature (under the discrete logarithm
  assumption on $\Gbn$).
\end{proposition}

% ── 8.3.1 ─────────────────────────────────────────────────────────────────────
\subsubsection{Robustness}
\label{sec:robustness}

The GOLDEN DKG provides \emph{verifiability} (any invalid dealing is detected
and rejected) but does \textbf{not} provide \emph{robustness}---the guarantee
that the protocol completes successfully even in the presence of malicious
participants.

Specifically, a single malicious dealer can prevent completion by broadcasting
an invalid dealing that every honest verifier rejects. Since
$\mathsf{Complete}$ requires exactly $n - 1$ verified peer dealings
(one from each other participant), a missing or invalid dealing causes all
honest participants to abort.

\begin{remark}[Robustness vs.\ Complaint Mechanisms]
  Achieving robustness typically requires a \emph{complaint round}
  (as in GJKR~\cite{gjkr1999}), where participants identify and exclude
  misbehaving dealers. GOLDEN's non-interactive design explicitly trades
  robustness for a single-round protocol: adding a complaint mechanism would
  require at least one additional interactive round, negating the main
  advantage of the construction.

  For applications requiring robustness, the recommended approach is to run
  GOLDEN DKG optimistically and fall back to a multi-round protocol (e.g.,
  GJKR or Pedersen DKG with complaints) if any dealing fails verification.
  The application layer should verify the number of valid dealings and decide
  whether to proceed (if $\ge t$ honest dealers produced valid dealings, the
  secret can still be reconstructed from a subset) or to abort and re-run.
\end{remark}

% ── 8.4 ──────────────────────────────────────────────────────────────────────
\subsection{Secret Material Lifecycle}
\label{sec:secret-lifecycle}

\begin{center}
\renewcommand{\arraystretch}{1.3}
\begin{tabular}{@{}llp{6cm}@{}}
\toprule
\textbf{Material} & \textbf{Created} & \textbf{Erased} \\
\midrule
$\omega_d$, $a_1, \ldots, a_{t-1}$ & Dealing & After share generation \\
$\{s_j\}_{j \ne d}$ & Dealing & After encryption \\
$\pad_{d,j}$ & Dealing / Completion & Immediately after use \\
Schnorr nonce $\nu$ & Identity proof & After proof generation \\
$\sk$ as \texttt{big.Int} + \texttt{skBytes} & PLONK witness & Word-level zeroing + \texttt{runtime.KeepAlive} \\
Decrypted shares $\hat{s}_{d,i}$ & Completion & After accumulation \\
$\mathsf{DkgDealing.PrivateShare}$ & Dealing & Caller invokes \texttt{DkgDealing.Zero()} \\
$\mathsf{DkgDealing.DerivedPrivateShares}$ & Dealing & Caller invokes \texttt{DkgDealing.Zero()} \\
$\mathsf{DkgDealing.RandomMsg}$ & Dealing & Caller invokes \texttt{DkgDealing.Zero()} (byte-by-byte) \\
$\mathsf{padBytesCache}$ & Dealing / Completion & Byte-by-byte zeroing after derived share processing via explicit closure \\
$\mathsf{DkgOutput.SecretShare}$ & Completion & Caller invokes \texttt{DkgOutput.Zero()} \\
$\mathsf{DkgOutput.DerivedOutputs[*].SecretShare}$ & Completion & Caller invokes \texttt{DkgOutput.Zero()} \\
FROST nonces $d_i, e_i$ & Signing Round 1 & After Round 2 (\texttt{defer}) \\
\midrule
\multicolumn{3}{@{}l}{\textit{Piggyback nonce preprocessing (Section~\ref{sec:piggyback})}} \\
Pending nonce $(d^{\mathrm{next}}, e^{\mathrm{next}})$ & $\mathsf{PiggybackSign}$ / $\mathsf{BootstrapRound1}$ & Zeroed when consumed in next session (step~16 of Algorithm~\ref{alg:piggyback-sign}) \\
Cloned nonce $(d^{\mathrm{clone}}, e^{\mathrm{clone}})$ & $\mathsf{PiggybackSign}$ step~15 & Zeroed by $\mathsf{SignRound2}$'s \texttt{defer} \\
\bottomrule
\end{tabular}
\end{center}

\begin{remark}[Go Memory Model Limitations]
  \label{rem:go-zeroing}
  All zeroing operations in the table above are \emph{best-effort} under the
  Go memory model. Three fundamental limitations apply:

  \begin{enumerate}[nosep]
    \item \textbf{GC compaction.} Go's garbage collector may copy heap objects
      during compaction, leaving stale copies of secret material in previously
      occupied memory. The \texttt{runtime.KeepAlive} calls prevent collection
      \emph{during} zeroing but cannot prevent prior copies.
    \item \textbf{Dead-store elimination.} The Go compiler may optimize away
      writes to variables that are not subsequently read.
      \texttt{runtime.KeepAlive} is the standard countermeasure, but it is not
      a formal guarantee against all compiler transformations.
    \item \textbf{Stack spills.} Local scalar variables may be spilled to the
      stack and not overwritten when the function returns.
  \end{enumerate}

  \noindent
  For applications with stringent secret erasure requirements (e.g., FIPS
  140-3), secrets should be allocated in non-GC-managed memory (e.g., via
  \texttt{mmap} with \texttt{MADV\_DONTDUMP} and explicit \texttt{mlock}),
  or the protocol should be implemented in a language with deterministic
  memory management (C, Rust). The zeroing in the current Go implementation
  represents the strongest practical mitigation available without
  \texttt{cgo} or \texttt{unsafe}.
\end{remark}

% ══════════════════════════════════════════════════════════════════════════════
\section{Protocol Data Flow}
\label{sec:dataflow}

Figure~\ref{fig:dataflow} illustrates the complete data flow from key
generation through DKG to FROST signing.

\begin{figure}[htbp]
\centering
\begin{tikzpicture}[
  box/.style={draw, rounded corners, minimum width=3cm, minimum height=0.8cm,
    text centered, font=\small},
  arrow/.style={-{Stealth[length=3mm]}, thick},
  dashedarrow/.style={-{Stealth[length=3mm]}, thick, dashed},
  group/.style={draw, dashed, rounded corners, inner sep=8pt,
    fill opacity=0.05, text opacity=1},
  node distance=0.8cm
]

% BJJ column
\node[box, fill=blue!10] (bjjkey) {$\sk_i \getsr \Fl$};
\node[box, fill=blue!10, below=of bjjkey] (bjjpk) {$\pk_i = [\sk_i]\genbjj$};
\node[box, fill=blue!10, below=of bjjpk] (dh) {$S = [\sk_d]\pk_j$};
\node[box, fill=blue!10, below=of dh] (xext) {$\dhx = S.X \in \Fr$};
\node[box, fill=blue!10, below=of xext] (h2c) {$H_1, H_2 \leftarrow \hashcurve(\mathsf{sessData})$};
\node[box, fill=blue!10, below=of h2c] (vrf) {$P_k = [\dhx \bmod \ell] H_k$};

% BN254 column (shifted right)
\node[box, fill=red!10, right=3.5cm of bjjkey] (bn254key) {$\omega_d \getsr \Fr$};
\node[box, fill=red!10, below=of bn254key] (poly) {$f_d(x) = \omega_d + \cdots$};
\node[box, fill=red!10, below=of poly] (vsscom) {$A_k = [a_k]\genbn$};
\node[box, fill=red!10, below=of vsscom] (shares) {$s_j = f_d(j) \in \Fr$};

% Cross-field
\node[box, fill=purple!10, below=3cm of shares] (lhl)
  {$\pad = P_1.X + \alpha \cdot P_2.X$};
\node[box, fill=red!10, below=of lhl] (encrypt) {$z = \pad + s_j$};
\node[box, fill=red!10, below=of encrypt] (rcommit) {$R = [\pad]\genbn$};

% ZK proof (center)
\node[box, fill=green!10, below=0.8cm of rcommit]
  (zkproof) {$\PLONK.\Prove(\mathsf{EVRFCircuit})$};

% Completion
\node[box, fill=orange!10, below=1cm of zkproof]
  (complete) {$\mathsf{secretShare}_i = \sum_d f_d(i)$};
\node[box, fill=orange!10, below=of complete]
  (groupkey) {$\mathsf{GroupKey} = \sum_d [\omega_d]\genbn$};

% FROST
\node[box, fill=yellow!20, below=of groupkey]
  (frost) {$\FROST.\mathsf{Sign}(m, \mathsf{KeyShare})$};
\node[box, fill=yellow!20, below=of frost]
  (sig) {$\sigma = (R, z)$};

% Piggyback
\node[box, fill=yellow!30, below=of sig]
  (piggyback) {$\mathsf{PiggybackSign}$: $(z_i, D_i^{\mathrm{next}}, E_i^{\mathrm{next}})$};

% Arrows
\draw[arrow] (bjjkey) -- (bjjpk);
\draw[arrow] (bjjpk) -- (dh);
\draw[arrow] (dh) -- (xext);
\draw[arrow] (xext) -- (h2c);
\draw[arrow] (h2c) -- (vrf);
\draw[arrow] (vrf) |- (lhl);

\draw[arrow] (bn254key) -- (poly);
\draw[arrow] (poly) -- (vsscom);
\draw[arrow] (poly) -- (shares);
\draw[arrow] (shares) |- (encrypt);

\draw[arrow] (lhl) -- (encrypt);
\draw[arrow] (lhl) -- (rcommit);
\draw[dashedarrow] (rcommit) -- (zkproof);
\draw[arrow] (encrypt) -- ++(0,-0.5) -| (complete);
\draw[arrow] (vsscom) -- ++(2,0) |- (groupkey);
\draw[arrow] (complete) -- (groupkey);
\draw[arrow] (groupkey) -- (frost);
\draw[arrow] (frost) -- (sig);
\draw[arrow] (sig) -- (piggyback);
\draw[arrow] (piggyback.east) -- ++(1.2,0) |- node[right, font=\scriptsize,
  pos=0.25, text width=2.2cm, align=left] {steady-state\\(1 round)} (frost.east);

% Group labels
\node[above=0.1cm of bjjkey, font=\footnotesize\bfseries, blue!70] {BJJ ($\Gbj$)};
\node[above=0.1cm of bn254key, font=\footnotesize\bfseries, red!70] {BN254 G1 ($\Gbn$)};

\end{tikzpicture}
\caption{Data flow from key generation through eVRF pad derivation,
  share encryption, DKG completion, and FROST signing. Blue boxes
  operate on BJJ, red boxes on BN254~G1, purple on the cross-field
  boundary, green on the PLONK proof system, orange on aggregation,
  and yellow on FROST signing.  The piggyback loop (bottom) shows the
  steady-state 1-round path after bootstrap
  (Section~\ref{sec:piggyback}).}
\label{fig:dataflow}
\end{figure}

\begin{figure}[htbp]
\centering
\begin{tikzpicture}[
  party/.style={draw, minimum width=2.2cm, minimum height=0.6cm,
    font=\small\bfseries, rounded corners},
  msg/.style={-{Stealth[length=2.5mm]}, thick},
  note/.style={font=\footnotesize, text width=3cm, align=center},
  timeline/.style={thick, dashed, gray},
  node distance=0.5cm
]

% Party headers
\node[party, fill=blue!15] (p1) at (0,0) {Party 1};
\node[party, fill=red!15] (p2) at (5,0) {Party 2};
\node[party, fill=green!15] (p3) at (10,0) {Party 3};

% Timelines
\draw[timeline] (0,-0.5) -- (0,-12.0);
\draw[timeline] (5,-0.5) -- (5,-12.0);
\draw[timeline] (10,-0.5) -- (10,-12.0);

% Phase 1: CreateDealing
\node[font=\footnotesize\itshape, gray] at (-2.2,-1.0) {CreateDealing};
\node[note] at (0,-1.0) {\footnotesize $\mathsf{msg}_1$};
\node[note] at (5,-1.0) {\footnotesize $\mathsf{msg}_2$};
\node[note] at (10,-1.0) {\footnotesize $\mathsf{msg}_3$};

% Broadcast arrows (Round 0)
\draw[msg] (0,-1.5) -- node[above, font=\scriptsize] {$\mathsf{msg}_1$} (5,-2.2);
\draw[msg] (0,-1.5) -- node[above, font=\scriptsize, pos=0.3] {} (10,-2.5);
\draw[msg] (5,-1.5) -- node[above, font=\scriptsize] {$\mathsf{msg}_2$} (0,-2.2);
\draw[msg] (5,-1.5) -- node[above, font=\scriptsize] {} (10,-2.2);
\draw[msg] (10,-1.5) -- node[above, font=\scriptsize, pos=0.7] {$\mathsf{msg}_3$} (0,-2.5);
\draw[msg] (10,-1.5) -- node[above, font=\scriptsize] {} (5,-2.5);

% Phase label
\draw[thick, dotted] (-3,-3.0) -- (12,-3.0);
\node[font=\footnotesize\itshape, gray] at (-2.2,-3.5) {VerifyDealing};

% Verify arrows (self-loops)
\draw[msg, blue!60] (0,-3.3) -- (0,-4.5)
  node[midway, left, font=\scriptsize, text width=2cm, align=right]
  {verify $\mathsf{msg}_2$, $\mathsf{msg}_3$};
\draw[msg, red!60] (5,-3.3) -- (5,-4.5)
  node[midway, left, font=\scriptsize, text width=2cm, align=right]
  {verify $\mathsf{msg}_1$, $\mathsf{msg}_3$};
\draw[msg, green!60] (10,-3.3) -- (10,-4.5)
  node[midway, right, font=\scriptsize, text width=2cm, align=left]
  {verify $\mathsf{msg}_1$, $\mathsf{msg}_2$};

% Phase label
\draw[thick, dotted] (-3,-5.0) -- (12,-5.0);
\node[font=\footnotesize\itshape, gray] at (-2.2,-5.5) {Complete};

% Complete
\draw[msg, blue!60] (0,-5.3) -- (0,-6.5)
  node[midway, left, font=\scriptsize, text width=2.5cm, align=right]
  {aggregate shares, compute group key};
\draw[msg, red!60] (5,-5.3) -- (5,-6.5)
  node[midway, left, font=\scriptsize] {};
\draw[msg, green!60] (10,-5.3) -- (10,-6.5)
  node[midway, right, font=\scriptsize] {};

% Output
\node[font=\scriptsize, blue!70] at (0,-7.0) {$(\mathsf{GK}, \mathsf{sk}_1, \ldots)$};
\node[font=\scriptsize, red!70] at (5,-7.0) {$(\mathsf{GK}, \mathsf{sk}_2, \ldots)$};
\node[font=\scriptsize, green!70] at (10,-7.0) {$(\mathsf{GK}, \mathsf{sk}_3, \ldots)$};

% Phase label
\draw[thick, dotted] (-3,-7.5) -- (12,-7.5);
\node[font=\footnotesize\itshape, gray] at (-2.2,-8.0) {FROST Sign};

% FROST
\node[font=\scriptsize] at (5,-8.3) {$t$-of-$n$ threshold signing};
\node[font=\scriptsize] at (5,-8.8) {$\sigma = (R, z)$};

% Phase label
\draw[thick, dotted] (-3,-9.3) -- (12,-9.3);
\node[font=\footnotesize\itshape, gray] at (-2.2,-9.8) {Piggyback};

% Piggyback bootstrap
\node[font=\scriptsize] at (5,-10.1) {Bootstrap: 2-round FROST (one-time)};

% Piggyback steady-state
\draw[msg, orange!80] (0,-10.5) -- node[above, font=\scriptsize] {$(z_1, D_1^{\mathrm{next}}, E_1^{\mathrm{next}})$} (5,-11.0);
\draw[msg, orange!80] (5,-10.5) -- node[above, font=\scriptsize] {$(z_2, D_2^{\mathrm{next}}, E_2^{\mathrm{next}})$} (10,-11.0);
\draw[msg, orange!80] (10,-10.5) -- node[above, font=\scriptsize, pos=0.7] {$(z_3, \ldots)$} (0,-11.0);

\node[font=\scriptsize] at (5,-11.5) {Steady-state: 1 round per session};

\end{tikzpicture}
\caption{Protocol sequence diagram for a 3-party GOLDEN DKG ($n = 3$). Each
  party broadcasts a \textsf{Round0Msg}, all parties verify all received
  dealings, then each party completes independently. The resulting key shares
  are used for FROST threshold signing.  After a 2-round bootstrap, piggyback
  nonce preprocessing (Section~\ref{sec:piggyback}) enables 1-round
  steady-state signing.}
\label{fig:sequence}
\end{figure}

% ══════════════════════════════════════════════════════════════════════════════
\section{Implementation Notes}
\label{sec:impl}

\subsection{Scalar Representations}

\begin{center}
\renewcommand{\arraystretch}{1.3}
\begin{tabular}{@{}llp{6cm}@{}}
\toprule
\textbf{Group} & \textbf{Scalar Type} & \textbf{Internal Representation} \\
\midrule
BJJ & \texttt{big.Int} & Arbitrary-precision integer, arithmetic mod $\ell$ \\
BN254 G1 & \texttt{fr.Element} & Montgomery form, $[4]\texttt{uint64}$ \\
\bottomrule
\end{tabular}
\end{center}

The Montgomery representation of BN254 scalars stores $a \cdot 2^{256} \bmod r$
rather than $a$ directly, enabling efficient modular multiplication via
Montgomery reduction. Serialization to bytes (for hashing, transmission, or
cross-field conversion) always produces the canonical integer representation.

\subsection{Point Representations}

\begin{center}
\renewcommand{\arraystretch}{1.3}
\begin{tabular}{@{}lll@{}}
\toprule
\textbf{Group} & \textbf{Affine} & \textbf{Intermediate} \\
\midrule
BJJ & \texttt{twistededwards.PointAffine} $(x, y)$ & Same (no projective) \\
BN254 G1 & \texttt{bn254.G1Affine} $(X, Y)$ & \texttt{G1Jac} for additions \\
\bottomrule
\end{tabular}
\end{center}

BN254 G1 uses Jacobian coordinates $(X:Y:Z)$ internally for point addition
(avoiding field inversions), converting to affine only for serialization and
equality checks.

\subsection{Circuit Constraint Count}

The eVRF circuit compiles to approximately $238{,}000$ constraints. The
dominant costs are:

\begin{center}
\renewcommand{\arraystretch}{1.3}
\begin{tabular}{@{}lr@{}}
\toprule
\textbf{Operation} & \textbf{Approx.\ Constraints} \\
\midrule
3 BJJ ScalarMul (native, $\sim$251-bit) & $3 \times 2{,}000 \approx 6{,}000$ \\
254-bit decomposition & $\sim 500$ \\
1 BN254 G1 ScalarMulBase (emulated) & $\sim 200{,}000$ \\
4 AssertIsOnCurve (BJJ, native) & $\sim 50$ \\
1 BN254 G1 AssertIsEqual (emulated) & $\sim 30{,}000$ \\
\midrule
\textbf{Total} & $\sim 238{,}000$ \\
\bottomrule
\end{tabular}
\end{center}

The emulated BN254 G1 operations dominate because each $\Fp'$ multiplication
requires multiple native $\Fr$ multiplications to simulate arithmetic in a
non-native field.

\subsection{Proof Size and Timing}

\begin{itemize}[nosep]
  \item PLONK proof size: $\sim 600$ bytes (bounded by \texttt{maxProofSize} $= 2048$).
    Proof size is validated against this bound during verification; oversized
    proofs are rejected before deserialization.
  \item Proof generation: $\sim 3$ seconds per proof (single-threaded).
  \item Proof verification: $\sim 10$ milliseconds.
  \item Per-dealing overhead: $n - 1$ proofs generated (one per recipient).
\end{itemize}

\subsection{CurveSuite Abstraction}
\label{sec:curvesuite}

The protocol logic is parameterized by a \texttt{CurveSuite} interface that
abstracts all curve-pair-specific operations. Implementing this interface for
a new curve pair (e.g., BLS12-381/JubJub) enables the GOLDEN DKG to operate
on different curve combinations without modifying the protocol code.

The interface requires the following methods:

\begin{center}
\renewcommand{\arraystretch}{1.3}
\begin{tabular}{@{}lp{8cm}@{}}
\toprule
\textbf{Method} & \textbf{Description} \\
\midrule
\texttt{InnerGroup() Group} & Returns the DH/identity curve group (e.g., BJJ). \\
\texttt{OuterGroup() Group} & Returns the commitment/VSS curve group (e.g., BN254 G1). \\
\texttt{H1(data ...[]byte) (Point, error)} & Hash-to-curve on the inner curve, domain 1. Result must be in the prime-order subgroup. \\
\texttt{H2(data ...[]byte) (Point, error)} & Hash-to-curve on the inner curve, domain 2. Independent from H1 (different domain separator). \\
\texttt{ExtractXAsOuterScalar(Point) (Scalar, error)} & Extract the x-coordinate of an inner-curve point and interpret it as an outer-group scalar. \\
\texttt{OuterToInnerScalar(Scalar) (Scalar, error)} & Convert an outer-group scalar to an inner-group scalar (reduction mod inner-group order). \\
\texttt{GenerateEVRFProof(...) ([]byte, error)} & Generate a ZK proof for eVRF pad derivation. Parameters: dealer SK, dealer PK, recipient PK, session data, $\alpha$, pad result. \\
\texttt{VerifyEVRFProof(...) error} & Verify a ZK proof for eVRF pad derivation. Parameters: dealer PK, recipient PK, session data, $\alpha$, R commitment, proof bytes. \\
\bottomrule
\end{tabular}
\end{center}

A new curve pair must satisfy the field relationship~\eqref{eq:fieldrel}
(inner base field $=$ outer scalar field) to enable native in-circuit
arithmetic. The H1/H2 functions must produce independent outputs for identical
inputs (via distinct domain separators) and return points in the prime-order
subgroup of the inner curve.

\subsection{Code Step Numbering}
\label{sec:step-numbering}

\begin{remark}
  The implementation's internal step numbering (visible in code comments) starts
  from the first cryptographic operation, after input validation. This creates
  an offset of approximately 1 from the Algorithm line numbers in this document.
  For example, ``Step~1'' in the code (sampling $\omega$) corresponds to
  Algorithm~\ref{alg:dealing}, line~3.
\end{remark}

\subsection{Error Sentinel Catalog}
\label{sec:errors}

The implementation defines 28~sentinel errors for structured error handling.
The following table documents each error and its trigger condition:

\begin{center}
\renewcommand{\arraystretch}{1.15}
\footnotesize
\begin{tabular}{@{}lp{7.5cm}@{}}
\toprule
\textbf{Error} & \textbf{Condition} \\
\midrule
\texttt{ErrInvalidNodeID} & Node ID not in $[1, \mathtt{MaxNodeID}]$ \\
\texttt{ErrInvalidConfig} & $T < 2$ or $N < T$ \\
\texttt{ErrDuplicateNodeID} & Duplicate participant ID detected \\
\texttt{ErrPeerCountMismatch} & Peer count $\ne N - 1$ \\
\texttt{ErrSessionIDMismatch} & Dealing session ID $\ne$ config session ID \\
\texttt{ErrIdentityPoint} & Point is the identity (degenerate) \\
\texttt{ErrCiphertextVerification} & $[z]\genbn \ne R + \VSS(j)$ \\
\texttt{ErrDegenerateDH} & DH shared secret is identity point \\
\texttt{ErrIdentityProofFailed} & Schnorr PoK verification failed \\
\texttt{ErrEVRFProofFailed} & PLONK eVRF proof verification failed \\
\texttt{ErrHashToCurveFailed} & Hash-to-curve exhausted 256 attempts \\
\texttt{ErrInvalidVSSLength} & $|\mathbf{A}| \ne T$ \\
\texttt{ErrMissingCiphertext} & Expected recipient has no ciphertext \\
\texttt{ErrMissingEVRFProof} & Expected recipient has no eVRF proof \\
\texttt{ErrDealerSelfCiphertext} & Dealer included ciphertext for self \\
\texttt{ErrIdentityRCommitment} & R commitment is identity point \\
\texttt{ErrDuplicateDealing} & Duplicate dealing from same dealer \\
\texttt{ErrUnknownDealer} & Dealing from unknown participant \\
\texttt{ErrSelfDealing} & Received own dealing in peer dealings \\
\texttt{ErrProofTooLarge} & Proof exceeds \texttt{maxProofSize} \\
\texttt{ErrNodeIDExceedsN} & Node ID $> N$ \\
\texttt{ErrExtraCiphertexts} & Extra ciphertexts beyond expected set \\
\texttt{ErrExtraEVRFProofs} & Extra eVRF proofs beyond expected set \\
\texttt{ErrNilIdentityProof} & Nil identity proof in dealing \\
\texttt{ErrNilCiphertext} & Nil ciphertext entry in dealing \\
\texttt{ErrNilSuite} & Nil CurveSuite passed to function \\
\texttt{ErrDerivedCurveCountMismatch} & Derived curve count $\ne$ config \\
\texttt{ErrDerivedCiphertextVerification} & Derived ciphertext check failed \\
\bottomrule
\end{tabular}
\end{center}

% ══════════════════════════════════════════════════════════════════════════════
\section{Comparison with Standard FROST DKG}
\label{sec:comparison}

\begin{center}
\renewcommand{\arraystretch}{1.3}
\begin{tabular}{@{}llll@{}}
\toprule
\textbf{Property} & \textbf{FROST DKG (Pedersen)} & \textbf{GJKR}~\cite{gjkr1999} & \textbf{GOLDEN DKG} \\
\midrule
Rounds & 2 (bcast + private) & 2 (bcast + private) & 1 (broadcast only) \\
Secure channels & Required & Required & Not required \\
Share encryption & Secure channel & Secure channel & eVRF-derived pads \\
ZK proofs & None & None & PLONK per recipient \\
Adaptive security & No & Yes (erasure model) & No \\
Identity curve & Same as signing & Same as signing & Separate (BJJ) \\
Signing curve & Single curve & Single curve & BN254 G1 \\
Msg size / dealer & $\bigO(n)$ & $\bigO(n)$ & $\bigO(n)$ ct + proofs \\
\bottomrule
\end{tabular}
\end{center}

\medskip
\noindent\textbf{Cost discussion.}
Each dealer generates $n - 1$ PLONK proofs, giving $n(n - 1) = \bigO(n^2)$
total proofs across all dealers. Each verifier checks all $n - 1$ proofs in
each of the $n - 1$ peer dealings, for $\bigO(n^2)$ verifications per party.
At ${\sim}3$~seconds per proof and ${\sim}10$~ms per verification, the protocol
is practical for $n \le 20$ (about 60 proofs per dealer, ${\sim}3$~minutes
generation time with single-threaded proving). Parallelizing proof generation
across cores or machines can extend this to larger groups.

The key trade-off is that GOLDEN eliminates the need for secure point-to-point
channels at the cost of PLONK proof generation per recipient. For applications
where establishing secure channels is expensive or impractical (e.g.,
blockchain-based DKG with broadcast-only communication), the GOLDEN approach is
advantageous.

% ══════════════════════════════════════════════════════════════════════════════
\section{Conclusion}
\label{sec:conclusion}

We have presented a complete specification of the GOLDEN non-interactive DKG
protocol adapted to BN254/Baby~Jubjub. The construction exploits the algebraic
identity between BJJ's base field and BN254's scalar field to build an
efficient eVRF pad derivation mechanism, proved correct via gnark PLONK
zero-knowledge proofs. The resulting key shares integrate seamlessly with the
FROST threshold Schnorr signing protocol.

The implementation provides the following guarantees:
\begin{itemize}[nosep]
  \item Non-interactive dealing via a single broadcast message per participant.
  \item Verifiable share encryption via Feldman VSS consistency checks and eVRF proofs.
  \item Session-bound identity proofs preventing cross-session replay.
  \item Information-theoretic secrecy of shares against $t-1$ corruptions.
  \item Direct FROST compatibility for threshold Schnorr signatures.
  \item Piggyback nonce preprocessing for 1-round steady-state signing after a one-time 2-round bootstrap (Section~\ref{sec:piggyback}).
  \item Multi-curve derived shares from a single DKG execution (Section~\ref{sec:derived}).
  \item Curve-agnostic protocol via the \texttt{CurveSuite} abstraction (Section~\ref{sec:curvesuite}).
  \item Systematic secure erasure of all intermediate secret material.
\end{itemize}

The production KZG SRS is sourced from the Aztec Ignition ceremony (176
participants); see Remark~\ref{rem:srs} for the integrity verification
procedure. The hash-to-curve construction should be evaluated for the target
threat model's timing requirements.

% ══════════════════════════════════════════════════════════════════════════════
\begin{thebibliography}{10}

\bibitem{golden2025}
  B.~B\"unz, Y.~Choi, and C.~Komlo.
  GOLDEN: Non-Interactive DKG.
  \emph{IACR ePrint Archive}, Report 2025/1924, 2025.

\bibitem{feldman1987}
  P.~Feldman.
  A practical scheme for non-interactive verifiable secret sharing.
  In \emph{FOCS}, pages 427--437, 1987.

\bibitem{shamir1979}
  A.~Shamir.
  How to share a secret.
  \emph{Communications of the ACM}, 22(11):612--613, 1979.

\bibitem{rfc9591}
  D.~Connolly, C.~Komlo, I.~Goldberg, and C.~A.~Wood.
  The Flexible Round-Optimized Schnorr Threshold (FROST) Protocol for
  Two-Round Schnorr Signatures.
  RFC~9591, Internet Engineering Task Force, June 2024.

\bibitem{gabizon2019}
  A.~Gabizon, Z.~J.~Williamson, and O.~Ciobotaru.
  PLONK: Permutations over Lagrange-bases for Oecumenical Noninteractive
  arguments of Knowledge.
  \emph{IACR ePrint Archive}, Report 2019/953, 2019.

\bibitem{gnark}
  ConsenSys.
  gnark: a fast zk-SNARK library.
  \url{https://github.com/Consensys/gnark}, 2020--2026.

\bibitem{lefkovits2017}
  B.~Whitehat, M.~Bellés, and J.~Baylina.
  Baby Jubjub Elliptic Curve.
  EIP-2494, Ethereum Improvement Proposals, 2020.

\bibitem{schwartz1980}
  J.~T.~Schwartz.
  Fast probabilistic algorithms for verification of polynomial identities.
  \emph{Journal of the ACM}, 27(4):701--717, 1980.

\bibitem{aztec-ignition}
  Aztec Protocol.
  Aztec Ignition: A Multi-Party Computation Ceremony for BN254.
  \url{https://github.com/AztecProtocol/ignition-verification}, 2019.

\bibitem{gjkr1999}
  R.~Gennaro, S.~Jarecki, H.~Krawczyk, and T.~Rabin.
  Secure distributed key generation for discrete-log based cryptosystems.
  In \emph{EUROCRYPT}, pages 295--310, 1999.

\bibitem{dodis2008}
  Y.~Dodis, R.~Ostrovsky, L.~Reyzin, and A.~Smith.
  Fuzzy extractors: How to generate strong keys from biometrics and other
  noisy data.
  \emph{SIAM Journal on Computing}, 38(1):97--139, 2008.

\bibitem{ill1989}
  R.~Impagliazzo, L.~Levin, and M.~Luby.
  Pseudo-random generation from one-way functions.
  In \emph{STOC}, pages 12--24, 1989.

\bibitem{kb2016}
  T.~Kim and R.~Barbulescu.
  Extended tower number field sieve: A new complexity for the medium prime
  case.
  In \emph{CRYPTO}, pages 543--571, 2016.

\end{thebibliography}

\end{document}
